% Lecture notes on the classification of homogeneous plane continua,
% following the formal (Lean 4) proof in this repository.
% Volume 7: The decomposable case.
%
% Compile with:  pdflatex main && pdflatex main   (two runs for the contents and references)
\documentclass[11pt,oneside]{book}

\usepackage[T1]{fontenc}
\usepackage{lmodern}
\usepackage[margin=1.15in]{geometry}
\usepackage{amsmath,amssymb,amsthm}
\usepackage{enumitem}
\usepackage{booktabs}
\usepackage{longtable}
\usepackage{tikz}
\usepackage[hidelinks]{hyperref}

% ---------------------------------------------------------------- theorem styles
\theoremstyle{plain}
\newtheorem{theorem}{Theorem}[chapter]
\newtheorem{lemma}[theorem]{Lemma}
\newtheorem{proposition}[theorem]{Proposition}
\newtheorem{corollary}[theorem]{Corollary}
\newtheorem{claim}[theorem]{Claim}
\theoremstyle{definition}
\newtheorem{definition}[theorem]{Definition}
\newtheorem{example}[theorem]{Example}
\newtheorem{exercise}[theorem]{Exercise}
\theoremstyle{remark}
\newtheorem{remark}[theorem]{Remark}

% ---------------------------------------------------------------- notation
\newcommand{\N}{\mathbb{N}}
\newcommand{\Z}{\mathbb{Z}}
\newcommand{\Q}{\mathbb{Q}}
\newcommand{\R}{\mathbb{R}}
\newcommand{\C}{\mathbb{C}}
\newcommand{\diam}{\operatorname{diam}}
\newcommand{\dist}{\operatorname{dist}}
\newcommand{\interior}{\operatorname{int}}
\newcommand{\cl}[1]{\overline{#1}}
\newcommand{\abs}[1]{\left\lvert #1\right\rvert}
\newcommand{\set}[1]{\left\{ #1 \right\}}
\newcommand{\nrm}[1]{\lVert #1\rVert}
\newcommand{\Log}{\operatorname{Log}}
\renewcommand{\Im}{\operatorname{Im}}
% the frontier of a set
\newcommand{\fr}{\operatorname{fr}}
% the element of a decomposition containing x (as in Volume 3)
\newcommand{\elt}[1]{[#1]}
% Jones's sets L_x and U_x
\newcommand{\Ls}[1]{L_{#1}}
\newcommand{\Us}[1]{U_{#1}}
% the component of the complement of K containing x (as in Volume 4)
\newcommand{\comp}[2]{\operatorname{C}_{#1}(#2)}
% the square annulus and the strip
\newcommand{\Ann}{\mathcal{A}}
\newcommand{\Strip}{\mathcal{S}}
% the right side, the left side and the not-right part of a crossing set of the strip
\newcommand{\Rc}{\operatorname{Rt}}
\newcommand{\Lc}{\operatorname{Lt}}
\newcommand{\NR}{\operatorname{NR}}
\newcommand{\Sec}{\operatorname{Sec}}
% the open arc
\newcommand{\oarc}[1]{\mathring{#1}}
% the unit interval and the circle
\newcommand{\I}{[0,1]}
\newcommand{\Sph}{S^1}

% Names of declarations in the Lean development (typewriter, underscores allowed).
\newcommand{\lean}[1]{\texttt{\detokenize{#1}}}
% A paragraph pointing to the formal counterpart of a result.
\newcommand{\formal}[1]{\par\smallskip\noindent{\small\textit{Formal counterpart:} #1}\par\smallskip}
% References to numbered results of the earlier volumes.
\newcommand{\vone}[1]{Volume~1, #1}
\newcommand{\vtwo}[1]{Volume~2, #1}
\newcommand{\vthree}[1]{Volume~3, #1}
\newcommand{\vfour}[1]{Volume~4, #1}
\newcommand{\vfive}[1]{Volume~5, #1}
\newcommand{\vsix}[1]{Volume~6, #1}

\setlist{itemsep=2pt,topsep=4pt}

\title{\textbf{Homogeneous Plane Continua}\\[6pt]
\Large Lecture notes following the formal proof\\[18pt]
\large Volume 7\\[4pt]
\Large The decomposable case:\\ aposyndesis, Jones's decomposition\\ and simple closed curves}
\author{}
\date{}

\begin{document}

\frontmatter
\maketitle
\tableofcontents

\mainmatter
% The introduction is Chapter 0, as in Volumes 1--6.
\setcounter{chapter}{-1}
%======================================================================
\chapter{Introduction}\label{ch:intro}
%======================================================================

\section{What this volume is about}

The classification of homogeneous plane continua (\vone{Theorem~0.1}) splits into an indecomposable and a decomposable case. Volume~6
treated the indecomposable case. This volume treats the decomposable case:

\begin{theorem}[Theorem~\ref{thm:decomposable}]
Let $M$ be a decomposable homogeneous plane continuum. Then either $M$ is a simple closed curve, or there is a continuous, open, surjective map
$f\colon M\to\Sph$ whose fibres are pseudo-arcs and are terminal in $M$ (every continuum meeting a fibre lies in it or contains it).
\end{theorem}

The proof follows two papers of F.~B.~Jones.
\begin{enumerate}
\item \textbf{Jones's decomposition} (1955, Theorem~1). For $x\in M$ let $\Ls x$ be the set of points $z$ such that every continuum containing a
neighbourhood of $z$ contains $x$. In a homogeneous continuum the sets $\Ls x$ partition $M$ into terminal continua; the decomposition is continuous;
and the quotient $M/G$ is a homogeneous \emph{aposyndetic} continuum. When $M$ lies in the plane, the elements do not separate the plane, and $M/G$
is again a plane continuum.
\item \textbf{Jones's theorem of 1949.} A nondegenerate homogeneous aposyndetic plane continuum is a simple closed curve. So $M/G$ is a circle
(Jones 1955, Theorem~2).
\item \textbf{The elements.} Then $M$ has exactly two complementary domains, so a decomposable element would separate the plane; hence the elements are
indecomposable, and by Volume~6 the nondegenerate ones are pseudo-arcs.
\end{enumerate}
The circle-like structure of $M$, and the fact that there is only one circle of pseudo-arcs up to homeomorphism, are the subject of Volume~8.

The volume has nine chapters after this introduction.

\begin{itemize}
\item \textbf{Chapter~\ref{ch:apos}: Aposyndesis.} Jones's 1948 paper: openness, the sets $\Ls y$ are continua, semi-local connectedness, decomposability.
\item \textbf{Chapter~\ref{ch:spread}: Homogeneity and aposyndesis.} Spreading (from the Effros theorem), the Lemma of Jones 1951, and Bing's theorem on non-cutting points.
\item \textbf{Chapter~\ref{ch:jones55}: Jones's aposyndetic decomposition.} Theorem~1 of 1955, continuity of the decomposition, the elements, non-separation of the plane;
Jones's 1951 theorem on hereditarily unicoherent continua.
\item \textbf{Chapter~\ref{ch:order}: Order and Menger's theorem.} Order of a point; a continuum all of whose points have order~$2$ is a circle.
\item \textbf{Chapter~\ref{ch:whyburn}: Whyburn's separating point--order theorem.} Only countably many local separating points fail to have order~$2$, by chains of separating points.
\item \textbf{Chapter~\ref{ch:sectors}: Continua crossing an annulus.} The sector theorem and the cyclic order, by lifting to a strip; crossing continua from a failure of local connectedness.
\item \textbf{Chapter~\ref{ch:boundaries}: Boundaries of complementary domains.} Whyburn's 1939 Theorem~14 and its Corollary~1: for a semi-locally connected plane continuum the
boundaries of complementary domains are locally connected, and simple closed curves if there are no cut points.
\item \textbf{Chapter~\ref{ch:jones49}: Homogeneous aposyndetic plane continua are circles.} Jones 1949, Theorems~1 and~2.
\item \textbf{Chapter~\ref{ch:decomposable}: The decomposable case.} The plane of elements, Jones 1955 Theorem~2 and its remarks, the elements are pseudo-arcs, the main theorem.
\end{itemize}

The main results are:

\begin{theorem}[Jones 1955, Theorem~1; Theorems~\ref{thm:jones55-b}, \ref{thm:jones55-c}, \ref{thm:jones55-usc}, \ref{thm:jones55-a}]
In a homogeneous continuum $X$ the sets $\Ls x$ partition $X$ into terminal subcontinua; the decomposition is continuous; and if $X$ is decomposable, the quotient is a
nondegenerate homogeneous aposyndetic continuum.
\end{theorem}

\begin{theorem}[Menger; Theorem~\ref{thm:menger}]
A nondegenerate continuum all of whose points have order $2$ is a simple closed curve.
\end{theorem}

\begin{theorem}[Whyburn; Theorem~\ref{thm:sep-order}, Corollary~\ref{cor:whyburn}]
In a separable locally compact metric space all but countably many local separating points have order $2$.
\end{theorem}

\begin{theorem}[Whyburn 1939; Theorems~\ref{thm:whyburn14} and~\ref{thm:whyburn-cor1}]
The boundary of each complementary domain of a semi-locally connected plane continuum is locally connected; if the continuum has no cut point (and is nondegenerate), the
boundary is a simple closed curve.
\end{theorem}

\begin{theorem}[Jones 1949; Theorem~\ref{thm:jones49-1}]
A nondegenerate homogeneous aposyndetic plane continuum is a simple closed curve.
\end{theorem}

\begin{theorem}[Jones 1955, Theorem~2; Theorem~\ref{thm:jones55-2}, Corollary~\ref{cor:jones55-remarks}]
If $M$ is a decomposable homogeneous plane continuum, $M/G$ is a simple closed curve, each element is a homogeneous indecomposable continuum not separating the plane, and $M$
has exactly two complementary domains.
\end{theorem}

\section{Where these results are used}

Theorem~\ref{thm:decomposable} is the decomposable case of the classification, assembled in Volume~8 together with the indecomposable case (\vsix{Theorem~7.4}) and the
theorem of Bing and Jones that circles of pseudo-arcs are unique. The map $f$ of Theorem~\ref{thm:decomposable} is the starting point of Volume~8.

\section{Where the formal proof differs from the literature}\label{sec:intro-diff}

\begin{enumerate}
\item \textbf{Jones 1955.} Parts (b) and (c) of Theorem~1 are proved from a consequence of the Effros theorem (spreading, Theorem~\ref{thm:spreading}) and the Lemma of Jones
1951, proved with a lower semicontinuous function in place of a transfinite sequence; Jones's twelve claims and his case argument are not needed (they are proved formally
nonetheless). Decomposability is not needed for (b) and (c). Aposyndesis of the quotient, which Jones states without argument, is proved from the continuity of the
decomposition. In the non-separation argument the disjointness of the domains $W_g$ is supplied (Theorem~\ref{thm:elt-not-separates}). The indecomposability of the elements is
derived from the two-domain theorem rather than from Jones's 1951 paper.
\item \textbf{Bing 1948.} Bing's Theorem~5 on non-cutting points is proved by a sequence of tight continua (Section~\ref{sec:bing5}).
\item \textbf{Menger and Whyburn.} Menger's theorem is proved in the special case needed, by a triod argument. Whyburn's separating point--order theorem is proved with chains
of separating points and a cut-wire argument, without his theory of cyclic elements and potential order.
\item \textbf{Moore 1922.} The two facts about crossing continua that Moore takes for granted (the sector theorem and the cyclic order) are proved by lifting to a strip, with
continuous logarithms only.
\item \textbf{Whyburn 1939.} The reduction to bounded domains by inversion has a gap and is avoided; Whyburn's arcs, his use of a theorem of Moore and his curve through a
point of order~$2$ are replaced by Janiszewski's theorem, the sector theorem and accumulating continua. Corollary~1, stated without proof, is proved with a $\theta$-curve lemma.
\item \textbf{Jones 1949.} Theorems~1 and~2 are false as printed for a one-point continuum. ``$M-(A+B)$ is not connected'' is proved with the $\theta$-curve lemma
(Lemma~\ref{lem:theta-sep}), and the two cases end with explicit local separating points.
\end{enumerate}

\section{Prerequisites and conventions}\label{sec:intro-conv}

The reader should know the material of Volume~3 (boundary bumping, irreducible continua, non-cut points, upper semicontinuous decompositions and their continuity at a
residual set of elements, the Effros theorem), the plane topology of Volume~4 (continuous logarithms, Janiszewski's theorem, Brouwer's theorem, invariance of plane separation)
and Volume~5 (arcs and simple closed curves, the Jordan curve theorem, Moore's theorem on decompositions of the plane). Chapter~\ref{ch:decomposable} uses the main theorem of
Volume~6.

\begin{itemize}
\item \textbf{Continua.} A continuum is compact, connected and metric; it may be a point unless ``nondegenerate'' is said.
\item \textbf{The plane.} As in Volumes~4 to~6, $\R^2$ carries the sup metric, so that balls are squares.
\item \textbf{Simple closed curves.} A simple closed curve is the union of two arcs with the same end points and no other common point; by \vfive{Proposition~1.13} it is
homeomorphic to the circle $\Sph$.
\end{itemize}

\paragraph{Formal counterparts.} As in the earlier volumes, many results are followed by a line ``Formal counterpart:'' naming the corresponding declaration of the
Lean development (all names begin with \lean{HPC.}). These lines may be skipped. Appendix~\ref{app:formal} collects them.

\paragraph{Exercises.} Each chapter ends with exercises. None of them is used later.

%======================================================================
\chapter{Aposyndesis}\label{ch:apos}
%======================================================================

Aposyndesis is F.~B.~Jones's name (1941) for a weak form of local connectedness: a continuum is
aposyndetic at $x$ with respect to $y$ if a continuum that misses $y$ contains a whole neighbourhood of
$x$. A decomposable continuum is aposyndetic at some point with respect to some other point, and the
points with respect to which a homogeneous continuum fails to be aposyndetic are the raw material of
Jones's decomposition (Chapter~\ref{ch:jones55}). This chapter collects what is needed from Jones's
1948 paper \emph{Concerning non-aposyndetic continua} and from Whyburn's 1939 paper
\emph{Semi-locally-connected sets}. Nothing here uses the plane or homogeneity.

\section{Conventions and definitions}

Throughout this chapter $X$ is a continuum, that is, a compact connected metric space; it may be a single
point. As in the formal development, Jones's continuum $M$ is the whole space $X$, and subsets of $X$
carry the subspace topology. A \emph{continuum in $X$} (or \emph{subcontinuum}) is a nonempty compact
connected subset. $X$ is \emph{decomposable} if $X=A\cup B$ with $A,B$ proper subcontinua.

\formal{The setting is a type \lean{X} with \lean{[MetricSpace X] [CompactSpace X] [ConnectedSpace X]};
continua in $X$ are \lean{HPC.Fold.Q78.IsContinuum}; decomposability is \lean{HPC.Jones.IsDecomposableSpace}.}

\begin{definition}\label{def:apos}
Let $x,y\in X$.
\begin{enumerate}[label=(\alph*)]
\item $X$ is \emph{aposyndetic at $x$ with respect to $y$} if there are a continuum $H$ and an open set $U$
with $x\in U\subseteq H\subseteq X\setminus\set y$.
\item $X$ is \emph{aposyndetic at $x$} if it is aposyndetic at $x$ with respect to every $y\neq x$, and
\emph{aposyndetic} if it is aposyndetic at every point. It is \emph{totally non-aposyndetic} if it is
aposyndetic at no point.
\item For $x\in X$ put
\[
\Us x=\set{z\in X : X\text{ is aposyndetic at }z\text{ with respect to }x},\qquad \Ls x=X\setminus\Us x .
\]
\end{enumerate}
\end{definition}

\formal{\lean{HPC.Jones.AposAt}, \lean{HPC.Jones.AposyndeticAt}, \lean{HPC.Jones.IsAposyndetic},
\lean{HPC.Jones.TotallyNonAposyndetic}; $\Us x$ is \lean{HPC.Jones.aposSet} (with
\lean{HPC.Jones.aposSet_eq}) and $\Ls x$ is \lean{HPC.Jones.Lset} (with \lean{HPC.Jones.Lset_compl}).}

Since a continuum $H$ containing a neighbourhood of $x$ contains $x$, $X$ is never aposyndetic at $x$
with respect to $x$. Hence $x\in\Ls x$, and $\Ls x$ is $x$ together with the points $z\ne x$ at which $X$
is not aposyndetic with respect to $x$. This is the set $L_x$ of Jones's 1955 paper, and $\Us x$ is the
set $U_x$ of his 1951 paper.

\begin{example}\label{ex:apos}
A Peano continuum (a locally connected continuum) is aposyndetic: for $x\ne y$, a connected open
neighbourhood $U$ of $x$ with $y\notin\cl U$ has a continuum $\cl U$ as its closure. An arc is
aposyndetic but has cut points. A nondegenerate indecomposable continuum is aposyndetic at no point with
respect to any other point (Theorem~\ref{thm:jones48-8} below). The product of two nondegenerate
continua is aposyndetic (Theorem~\ref{thm:jones48-7}).
\end{example}

\begin{definition}\label{def:slc}
$X$ is \emph{semi-locally connected at $x$} if for every open $U\ni x$ there is an open $V$ with
$x\in V\subseteq U$ such that $X\setminus U$ lies in the union of finitely many components of $X\setminus V$.
\end{definition}

\formal{\lean{HPC.Jones.SemiLocConnAt}.}

\begin{definition}\label{def:cuts}
A point $p$ \emph{cuts $x$ from $y$} if $x$, $y$, $p$ are distinct and every continuum containing $x$ and
$y$ contains $p$. A \emph{weak cut point} is a point that cuts some $x$ from some $y$. A \emph{cut point} of
a connected space $Y$ is a point $p$ with $Y\setminus\set p$ not connected.
\end{definition}

\formal{\lean{HPC.Jones.CutsPts}, \lean{HPC.Jones.IsWeakCutPoint}, \lean{HPC.Jones.IsStrongCutPoint}.}

\section{Openness and finite covers}

\begin{lemma}\label{lem:apos-basic}
\leavevmode
\begin{enumerate}[label=(\alph*)]
\item If $H$, $U$ witness aposyndesis at $x$ with respect to $y$, then $X$ is aposyndetic at every point of
$U$ with respect to every point outside $H$.
\item If $T\colon X\to X$ is a homeomorphism, $X$ is aposyndetic at $x$ with respect to $y$ if and only if
it is aposyndetic at $Tx$ with respect to $Ty$. Hence $T(\Us x)=\Us{Tx}$ and $T(\Ls x)=\Ls{Tx}$.
\end{enumerate}
\end{lemma}

\begin{proof}
(a) is immediate from the definition. (b) $T(H)$ is a continuum and $T(U)$ is open; apply the same to
$T^{-1}$.
\end{proof}

\formal{\lean{HPC.Jones.aposAt_of_mem}, \lean{HPC.Jones.aposAt_homeomorph_iff},
\lean{HPC.Jones.aposSet_homeomorph}, \lean{HPC.Jones.Lset_homeomorph}.}

\begin{theorem}[Jones 1948, Theorem~1]\label{thm:jones48-1}
The set of pairs $(x,y)\in X\times X$ such that $X$ is aposyndetic at $x$ with respect to $y$ is open in
$X\times X$. Consequently, for fixed $y$ the set $\Us y$ is open and $\Ls y$ is closed; and if
$x_i\to x$, $y_i\to y$ and $X$ is not aposyndetic at $x_i$ with respect to $y_i$, then it is not
aposyndetic at $x$ with respect to $y$.
\end{theorem}

\begin{proof}
Let $H$, $U$ witness aposyndesis at $x$ with respect to $y$. Since $H$ is closed and $y\notin H$, the set
$X\setminus H$ is an open neighbourhood of $y$. By Lemma~\ref{lem:apos-basic}(a) every pair in the open set
$U\times(X\setminus H)$ is a pair of aposyndesis.
\end{proof}

\formal{\lean{HPC.Jones.isOpen_setOf_aposAt}, \lean{HPC.Jones.isOpen_aposSet},
\lean{HPC.Jones.isClosed_Lset}, \lean{HPC.Jones.theorem1}. Jones states the sequential form with the
hypotheses $x\ne y$, $x_i\ne y_i$, which are not needed.}

\begin{lemma}[Heine--Borel for aposyndesis]\label{lem:heine-borel}
Let $K\subseteq X$ be compact.
\begin{enumerate}[label=(\alph*)]
\item If $X$ is aposyndetic at every point of $K$ with respect to $y$, there are finitely many continua
$H_1,\dots,H_k$ missing $y$ and open sets $U_i\subseteq H_i$ with $K\subseteq U_1\cup\dots\cup U_k$.
\item If $X$ is aposyndetic at $x$ with respect to every point of $K$, there are an open $U\ni x$ and, for
each $y\in K$, a continuum $H_y\supseteq U$ missing $y$.
\end{enumerate}
\end{lemma}

\begin{proof}
(a) Cover $K$ by the open sets $U$ of witnesses $(H,U)$ and take a finite subcover. (b) For $y\in K$ take a
witness $(H_y,U_y)$; by Lemma~\ref{lem:apos-basic}(a) the same witness works for all points of the open set
$X\setminus H_y\ni y$. Finitely many of these sets cover $K$; the intersection of the corresponding $U_y$
is the required $U$.
\end{proof}

\formal{\lean{HPC.Jones.exists_finite_apos_cover}, \lean{HPC.Jones.exists_nhd_uniform_apos}.}

The next lemma is the tool for proving aposyndesis without exhibiting a continuum.

\begin{lemma}[Aposyndesis from the frontier]\label{lem:apos-frontier}
Let $V\subseteq X$ be open, $b\in V$ and $y\notin\cl V$. If $X$ is aposyndetic at every point of the
frontier $\fr V$ with respect to $y$, then it is aposyndetic at $b$ with respect to $y$.
\end{lemma}

\begin{proof}
By Lemma~\ref{lem:heine-borel}(a) there are continua $H_1,\dots,H_k$ missing $y$ and open $U_i\subseteq H_i$
covering the compact set $\fr V$. Let $F=\cl V\cup H_1\cup\dots\cup H_k$, a closed set missing $y$.

Every point $z$ of $F$ lies in the component of $F$ of some $H_i$. This is clear if $z\in H_i$. If $z\in\cl
V$, then $\cl V\ne X$ (it misses $y$), and by boundary bumping (\vthree{Theorem~1.10}) the component of $z$
in $\cl V$ meets $\cl{X\setminus\cl V}$, hence meets $\fr V$, hence meets some $U_i\subseteq H_i$; so $z$
lies with $H_i$ in a connected subset of $F$.

So $F$ has at most $k$ components $C_1,\dots,C_k$, and they cover the open set $O=V\cup U_1\cup\dots\cup U_k$.
Let $C$ be the component of $b$. The components other than $C$ are closed and miss $b$; removing them from
$O$ leaves an open set containing $b$ and contained in $C$. And $C$ is a continuum (a component of a
closed set) missing $y$.
\end{proof}

\formal{\lean{HPC.Jones.aposAt_of_frontier}; the last step is \lean{HPC.Jones.exists_open_subset_component}.}

\section{The sets \texorpdfstring{$\Ls y$}{L\_y} are continua}

\begin{theorem}[Jones 1948, Theorem~3]\label{thm:jones48-3}
For every $y\in X$ the set $\Ls y$ is a continuum.
\end{theorem}

\begin{proof}
$\Ls y$ is closed by Theorem~\ref{thm:jones48-1} and contains $y$. Suppose $\Ls y=A\cup B$ with $A$, $B$ closed,
disjoint and nonempty, $y\in A$. Choose an open $V\supseteq B$ with $\cl V\cap A=\emptyset$. Then $\fr V$
misses $A$ (as $\fr V\subseteq\cl V$) and misses $B$ (as $B\subseteq V$ and $V$ is open), so $X$ is
aposyndetic at every point of $\fr V$ with respect to $y$. Also $y\in A$ is not in $\cl V$. By
Lemma~\ref{lem:apos-frontier}, $X$ is aposyndetic at every point of $B$ with respect to $y$, against
$B\subseteq\Ls y$.
\end{proof}

\formal{\lean{HPC.Jones.theorem3}, \lean{HPC.Jones.isContinuum_Lset}. Jones prints no proof.}

\section{Semi-local connectedness}

\begin{theorem}[Jones 1948, Theorem~0]\label{thm:jones48-0}
$X$ is semi-locally connected at $p$ if and only if $X$ is aposyndetic at every $x\ne p$ with respect
to $p$.
\end{theorem}

\begin{proof}
Suppose $X$ is semi-locally connected at $p$, and let $x\ne p$, $d=d(x,p)$. Apply the definition to
$U=B(p,d/2)$: there is an open $V$ with $p\in V\subseteq U$ such that $X\setminus U$ lies in finitely many
components of $X\setminus V$. The open set $O=X\setminus\bar B(p,d/2)$ contains $x$ and lies in $X\setminus U$,
hence in the union of finitely many components of the closed set $X\setminus V$. As in the last step of the
proof of Lemma~\ref{lem:apos-frontier}, the component $C$ of $x$ in $X\setminus V$ contains an open set
containing $x$. $C$ is a continuum missing $p\in V$.

Conversely, suppose the condition and let $U\ni p$ be open. By Lemma~\ref{lem:heine-borel}(a) for the
compact set $K=X\setminus U$, there are continua $H_1,\dots,H_k$ missing $p$ whose interiors cover $X\setminus U$.
Put $V=U\setminus(H_1\cup\dots\cup H_k)$, an open set containing $p$. Each $H_i$ is a connected subset of
$X\setminus V$, so lies in one component of $X\setminus V$, and these $k$ components cover $X\setminus U$.
\end{proof}

\formal{\lean{HPC.Jones.theorem0}. Jones only says that it ``may be established with the help of the
Heine--Borel theorem''.}

\begin{corollary}[Jones 1941, Theorem~4]\label{cor:jones41}
A continuum is aposyndetic if and only if it is semi-locally connected at each of its points.
\end{corollary}

\begin{proof}
Both conditions say that $X$ is aposyndetic at $x$ with respect to $p$ for all $x\ne p$.
\end{proof}

\formal{\lean{HPC.Jones.semiLocConn_of_isAposyndetic}, \lean{HPC.Jones.isAposyndetic_of_semiLocConn}.}

\begin{proposition}[Whyburn 1939, (6.2)]\label{prop:whyburn62}
Let $X$ be semi-locally connected at each point, $x\in X$, and $a\ne x$. If $b$ lies in the component of
$a$ in $X\setminus\set x$, then some continuum contains $a$ and $b$ and misses $x$.
\end{proposition}

\begin{proof}
Let $S$ be the set of points $z\ne x$ such that some continuum contains $a$ and $z$ and misses $x$. Then
$a\in S$. By Corollary~\ref{cor:jones41}, $X$ is aposyndetic at every $z\neq x$ with respect to $x$; let $H$
be a continuum missing $x$ and containing an open $W\ni z$. If $z\in S$, with continuum $K$, then every
point of $W$ is in $S$ (use $K\cup H$); so $S$ is open. If $z\ne x$ is a limit point of $S$, then $W$ contains a
point $z'\in S$ with continuum $K'$, and $K'\cup H$ shows $z\in S$; so $S$ is closed in $X\setminus\set x$. Thus $S$
is open and closed in $X\setminus\set x$ and contains the component of $a$.
\end{proof}

\formal{\lean{HPC.Jones.whyburn_6_2}; also Whyburn's (6.1) and (6.21), \lean{HPC.Jones.whyburn_6_1},
\lean{HPC.Jones.whyburn_6_21}. Whyburn argues with the components of $X\setminus V_\delta(x)$; the
formal proof goes through Theorem~\ref{thm:jones48-0}.}

\section{Decomposability}

\begin{theorem}[Jones 1948, Theorem~8]\label{thm:jones48-8}
$X$ is decomposable if and only if there are $x\ne y$ such that $X$ is aposyndetic at $x$ with respect
to $y$. Equivalently, $X$ is indecomposable if and only if $\Ls y=X$ for every $y$.
\end{theorem}

\begin{proof}
If $X=A\cup B$ with $A,B$ proper subcontinua, take $x\in X\setminus B$ and $y\in X\setminus A$. Then $A$ is a
continuum missing $y$ and containing the open set $X\setminus B\ni x$. In particular $x\ne y$.

Conversely, let $H$ be a continuum missing $y$ and containing an open $U\ni x$, with $x\ne y$. If $X\setminus H$ is
connected, then $X=H\cup\cl{X\setminus H}$, where $H\ne X$ (it misses $y$) and $\cl{X\setminus H}\neq X$ (it misses
the open set $U\subseteq H$). If $X\setminus H=A\cup B$ with $A,B$ separated and nonempty, then $H\cup A$ and $H\cup B$
are continua (\vthree{Lemma~1.5}), proper (they miss $B$, respectively $A$), and their union is $X$.
\end{proof}

\formal{\lean{HPC.Jones.theorem8}, \lean{HPC.Jones.theorem9}; the first direction is
\lean{HPC.Jones.aposAt_of_decomposition}.}

\begin{theorem}[Jones 1948, Theorem~7]\label{thm:jones48-7}
The product $Y\times Z$ of two nondegenerate continua is aposyndetic.
\end{theorem}

\begin{proof}
Let $a=(a_1,a_2)\ne b=(b_1,b_2)$. Suppose $a_1\ne b_1$; let $A=\bar B(a_1,d(a_1,b_1)/2)\subseteq Y$ and choose
$c\in Z$ with $c\ne b_2$. The set $H=(A\times Z)\cup(Y\times\set c)$ is a continuum (each slice $\set{y}\times
Z$ with $y\in A$ meets the connected set $Y\times\set c$), it misses $b$, and it contains the open set
$B(a_1,d(a_1,b_1)/2)\times Z\ni a$. The case $a_2\ne b_2$ is symmetric.
\end{proof}

\formal{\lean{HPC.Jones.theorem7}, with \lean{HPC.Jones.isContinuum_prod_cross}.}

\begin{remark}[The rest of Jones 1948]\label{rem:jones48-rest}
Jones's paper continues with weak cut points. Theorem~6 says: if $D$ is open and nonempty, $M_0$ closed and
disjoint from $D$, and for each $x\in D$ there is $y\in M_0$ with $X$ not aposyndetic at $x$ with respect to $y$,
then for each $z\notin M_0$ some $y_0\in M_0$ cuts some $x_0\in D$ from $z$. Theorems~11--19 and three corollaries
derive the existence and density of weak cut points; for instance, a totally non-aposyndetic continuum
contains a weak cut point, and if $X$ has no weak cut point, the points at which $X$ is aposyndetic and
semi-locally connected form a dense $G_\delta$ (Theorem~18). Only Theorem~18 reappears in these notes, in
Jones's Theorem~2 of 1949 (Corollary~\ref{cor:jones49-2}). Two statements are false as printed: Theorem~6
without the word ``nonempty'' (take $D=M_0=\emptyset$), and Theorem~10, ``$X$ is indecomposable iff it contains
an open set $D$ such that $X$ is not aposyndetic at any point of $D$ with respect to any other point'', whose
right-hand side always holds with $D=\emptyset$. Both are proved with $D$ nonempty.
\end{remark}

\formal{\lean{HPC.Jones.theorem6}, \lean{HPC.Jones.not_theorem6_printed}, \lean{HPC.Jones.theorem10},
\lean{HPC.Jones.theorem10_printed_iff}, \lean{HPC.Jones.not_theorem10_printed_unitInterval},
\lean{HPC.Jones.theorem11}, \lean{HPC.Jones.theorem12}, \lean{HPC.Jones.theorem13},
\lean{HPC.Jones.totallyNonApos_exists_weakCutPoint}, \lean{HPC.Jones.theorem18}
(in \lean{HPC/Jones/Basic.lean}, \lean{HPC/Jones/Theorem6.lean}, \lean{HPC/Jones/Theorem13.lean},
\lean{HPC/Jones/CutPoints.lean}).}

\section*{Exercises}

\begin{exercise}
Show that the $\sin(1/x)$-continuum is not aposyndetic, and determine the sets $\Ls y$ for all its points
$y$.
\end{exercise}

\begin{exercise}
Let $F$ be the harmonic fan: the union of the segments from $v=(0,0)$ to $(1,1/n)$, $n\ge1$, and to $(1,0)$.
Show that a point $p\neq v$ of the last segment lies in $\Ls v$, but $v\notin\Ls p$. Conclude that the sets $\Ls y$ of a
continuum need not be pairwise disjoint or equal. Chapter~\ref{ch:jones55} shows that in a homogeneous continuum they
are.
\end{exercise}

\begin{exercise}
Prove directly, without Theorem~\ref{thm:jones48-0}, that a locally connected continuum is semi-locally
connected at every point.
\end{exercise}

\begin{exercise}
Show that Lemma~\ref{lem:apos-frontier} fails without the hypothesis $y\notin\cl V$.
\end{exercise}

%======================================================================
\chapter{Homogeneity and aposyndesis}\label{ch:spread}
%======================================================================

This chapter proves three facts about a homogeneous continuum $X$. The first is a consequence of the
Effros theorem that the formal proof calls \emph{spreading}: a continuum with nonempty interior that
misses a point $x$ can be enlarged to a continuum that still misses $x$ and contains a neighbourhood of
any given one of its points. The second is the Lemma of Jones's 1951 paper: the sets $\Us x$ are never
properly nested. The third is Bing's theorem (1948) that some point of a nondegenerate continuum does
not cut any two open sets from each other; with homogeneity, no point does. These three facts replace
most of the long case analysis of Jones's 1955 paper.

Throughout, $X$ is a compact metric space; from Section~\ref{sec:bing5} on it is also connected.
\emph{Homogeneous} means that for all $x,y\in X$ some homeomorphism $T\colon X\to X$ has $Tx=y$.

\formal{\lean{HPC.Jones.IsHomogeneous}; it is the hypothesis of \lean{HPC.Effros.effros_homogeneous}
(\lean{HPC.Jones.isHomogeneous_iff_isPretransitive}).}

\section{Spreading}

\begin{theorem}[Spreading]\label{thm:spreading}
Let $X$ be a homogeneous compact metric space and $K\subseteq X$ a continuum with nonempty interior. If
$o\in K$ and $x\notin K$, then $X$ is aposyndetic at $o$ with respect to $x$.
\end{theorem}

\begin{proof}
Let $b$ be an interior point of $K$ and choose $\varepsilon_0>0$ with $B(b,\varepsilon_0)\subseteq K$. Let
\[
\varepsilon=\min\bigl(\varepsilon_0,\ \tfrac12\dist(x,K)\bigr)>0 ,
\]
and let $\mathcal S$ be the set of homeomorphisms $h$ of $X$ with $d(h(z),z)<\varepsilon$ for all $z$. Put
\[
Z=\cl{\bigcup_{h\in\mathcal S}\bigl(K\cup h(K)\bigr)} .
\]
\begin{itemize}
\item \emph{$Z$ is a continuum.} Each $K\cup h(K)$ is connected: $h(b)\in B(b,\varepsilon)\subseteq K$, so the
connected sets $K$ and $h(K)$ meet. All these sets contain $o$, so their union is connected, and so is its
closure. $Z$ is closed in the compact space $X$.
\item \emph{$x\notin Z$.} Every point of $h(K)$ is within $\varepsilon$ of $K$, so every point of $Z$ is at
distance $\le\varepsilon\le\frac12\dist(x,K)$ from $K$.
\item \emph{$Z$ contains a neighbourhood of $o$.} By the Effros theorem (\vthree{Theorem~7.15}) there is
$\delta>0$ such that every $y$ with $d(o,y)<\delta$ is $h(o)$ for some $h\in\mathcal S$; then $y\in h(K)\subseteq Z$. So
$B(o,\delta)\subseteq Z$.\qedhere
\end{itemize}
\end{proof}

\formal{\lean{HPC.Jones.aposAt_of_isContinuum_interior}, from \lean{HPC.Effros.effros_homogeneous}.}

\begin{corollary}[Non-aposyndesis is transitive]\label{cor:trans}
Let $X$ be homogeneous. If $X$ is not aposyndetic at $z$ with respect to $y$, nor at $y$ with respect to
$x$, then it is not aposyndetic at $z$ with respect to $x$. In terms of the sets of
Definition~\ref{def:apos}: if $z\in\Ls y$ and $y\in\Ls x$, then $z\in\Ls x$.
\end{corollary}

\begin{proof}
Suppose $K$ is a continuum missing $x$ that contains an open $V\ni z$. Then $y\in K$, for otherwise $(K,V)$
would make $X$ aposyndetic at $z$ with respect to $y$. Now $K$ has nonempty interior, contains $y$ and misses
$x$; by Theorem~\ref{thm:spreading}, $X$ is aposyndetic at $y$ with respect to $x$, a contradiction.
\end{proof}

\formal{\lean{HPC.Jones.not_aposAt_trans}.}

\section{The Lemma of Jones 1951}

\begin{lemma}\label{lem:open-subset-U}
For every closed $K\subseteq X$, the set $\set{w\in X : K\subseteq\Us w}$ is open.
\end{lemma}

\begin{proof}
Let $K\subseteq\Us w$, that is, $X$ is aposyndetic at every point of $K$ with respect to $w$. By
Lemma~\ref{lem:heine-borel}(a), finitely many continua $H_i$ missing $w$ have interiors covering $K$. For $w'$
in the open set $X\setminus\bigcup_iH_i$ the same continua show $K\subseteq\Us{w'}$.
\end{proof}

\formal{\lean{HPC.Jones.isOpen_setOf_subset_aposSet}.}

\begin{theorem}[Jones 1951, Lemma]\label{thm:jones51-lemma}
If $X$ is homogeneous, then $\Us y$ is not a proper subset of $\Us x$, for any $x,y\in X$.
\end{theorem}

\begin{proof}
Let $B_0,B_1,\dots$ be a countable basis of $X$ and define
\[
g(w)=\sum_{k\ge0}2^{-k}\,\chi_k(w),\qquad \chi_k(w)=\begin{cases}1&\text{if }\cl{B_k}\subseteq\Us w,\\ 0&\text{otherwise.}\end{cases}
\]
By Lemma~\ref{lem:open-subset-U} each $\chi_k$ is the indicator function of an open set, hence lower
semicontinuous, and so is the sum $g$. A lower semicontinuous function on a compact space attains its
minimum, say at $w_0$.

If $\Us a\subsetneq\Us b$, then $g(a)<g(b)$: every term of $g(a)$ is at most the corresponding term of $g(b)$;
and for $z\in\Us b\setminus\Us a$ there is $k$ with $z\in B_k\subseteq\cl{B_k}\subseteq\Us b$ (the set $\Us b$ is open,
Theorem~\ref{thm:jones48-1}), and then $\chi_k(b)=1$, $\chi_k(a)=0$ because $z\in\cl{B_k}\setminus\Us a$.

Now suppose $\Us y\subsetneq\Us x$ and let $T$ be a homeomorphism with $Tx=w_0$. By
Lemma~\ref{lem:apos-basic}(b), $\Us{Ty}=T(\Us y)\subsetneq T(\Us x)=\Us{w_0}$, so $g(Ty)<g(w_0)$, against the
minimality of $w_0$.
\end{proof}

\formal{\lean{HPC.Jones.jones51_lemma}, with the function $g$ from
\lean{HPC.Jones.exists_lsc_strictMono_aposSet}. Jones's proof runs a transfinite sequence of points
$x_0,x_1,\dots$ through the countable ordinals with $\Us{x_i}$ strictly decreasing, and contradicts second
countability; the lower semicontinuous function replaces the transfinite recursion.}

\section{Bing's theorem on non-cutting points}\label{sec:bing5}

From now on $X$ is a nondegenerate continuum. Bing's paper \emph{Some characterizations of arcs and simple
closed curves} (1948) contains the following theorem, which Jones's 1955 paper quotes as
``Corollary~2 of [Bing]''.

\begin{theorem}[Bing 1948, Theorem~5, point form]\label{thm:bing5}
For every $r\in X$ there is a point $p\ne r$ such that the union of the subcontinua of $X$ which contain $r$
and miss $p$ is dense in $X$.
\end{theorem}

The proof needs a definition and three lemmas. Call a subcontinuum $N$ of $X$ \emph{tight} if there are no
proper subcontinuum $K$ of $X$ and open set $O$ with $N\subseteq O\subseteq K$; that is, $N$ lies in the interior of no
proper subcontinuum.

\begin{lemma}\label{lem:tight-conn}
Let $N\subseteq E$ be subcontinua of $X$ with $N$ tight. Then $\cl{X\setminus E}\cup N$ is connected.
\end{lemma}

\begin{proof}
Let $G=\cl{X\setminus E}\cup N$ and suppose $G=A\cup B$ with $A$, $B$ closed, disjoint and nonempty; as $N$ is
connected we may take $N\subseteq A$. Choose disjoint open sets $u\supseteq A$, $v\supseteq B$ and put $Q=A\cup E$, a closed set.

\emph{$Q$ is connected.} Otherwise $Q=a\cup b$ with $a$, $b$ closed, disjoint and nonempty, and $E\subseteq a$; then
$b\subseteq A\setminus E$. Every point of $X\setminus Q$ lies in $X\setminus E\subseteq G$ but not in $A$, hence in $B$. So
$X=b\cup(a\cup B)$, where $b$ and $a\cup B$ are closed, nonempty and disjoint (because $b\subseteq A$). This
contradicts the connectedness of $X$.

\emph{$Q\ne X$.} $B$ is nonempty and contained in $\cl{X\setminus E}$ (as $B\cap N=\emptyset$), so the open set $v\supseteq B$
contains a point $z\notin E$; and $z\notin A$ since $A\subseteq u$. So $z\notin Q$.

\emph{$u\subseteq Q$.} A point of $u$ outside $E$ lies in $G=A\cup B$ but not in $B\subseteq v$, so it lies in $A$.

Thus $Q$ is a proper subcontinuum containing the open set $u\supseteq N$, against the tightness of $N$.
\end{proof}

\formal{\lean{HPC.Plane.Arcs48.isPreconnected_closure_diff_union} (Bing: ``$F_i+N$ is a continuum'').}

\begin{lemma}\label{lem:tight-piece}
Let $N$ be a tight subcontinuum and $D$ a nonempty open set with $D\cap N=\emptyset$. There is a closed set $Z\subseteq X$
with empty interior such that for every $p\in X\setminus(N\cup Z)$ some continuum $G\supseteq N$ misses $p$ and meets $\cl
D$.
\end{lemma}

\begin{proof}
Fix $n_0\in N$. The closed set $X\setminus D$ is proper and contains $n_0$. By boundary bumping
(\vthree{Theorem~1.10}) the component $K$ of $n_0$ in $X\setminus D$ meets $\cl D$. Let $E=K\cup N$, a continuum
contained in $X\setminus D$, and $Z=E\cap\cl{X\setminus E}$. A nonempty open set inside $Z$ would lie in $E$ and meet
$X\setminus E$ (being inside $\cl{X\setminus E}$); so $Z$ has empty interior. Let $p\notin N\cup Z$. If $p\notin E$,
take $G=E$. If $p\in E$, then $p\notin\cl{X\setminus E}$; take $G=\cl{X\setminus E}\cup N$, a continuum by
Lemma~\ref{lem:tight-conn}, which contains $D\subseteq X\setminus E$ and misses $p$.
\end{proof}

\formal{\lean{HPC.Plane.Arcs48.exists_nowhereDense_piece}.}

\begin{lemma}[The tight case]\label{lem:tight-case}
If $N$ is a tight proper subcontinuum, there is $p\in X\setminus N$ such that the union of the continua containing $N$
and missing $p$ is dense in $X$.
\end{lemma}

\begin{proof}
Let $b_0,b_1,\dots$ be a countable basis. For each $k$ with $b_k\ne\emptyset$ and $\cl{b_k}\cap N=\emptyset$ let $Z_k$ be the
set of Lemma~\ref{lem:tight-piece} for $D=b_k$. The sets $Z_k$ are closed with empty interior and $X\setminus N$ is a
nonempty open set, so by the Baire category theorem some $p\in X\setminus N$ lies in no $Z_k$. Let $x\in X$ and $W\ni x$
open. If $x\in N$, the continuum $N$ itself meets $W$. Otherwise choose $k$ with $x\in b_k\subseteq\cl{b_k}\subseteq W\setminus N$; the
continuum $G$ given by Lemma~\ref{lem:tight-piece} contains $N$, misses $p$ and meets $\cl{b_k}\subseteq W$.
\end{proof}

\formal{\lean{HPC.Plane.Arcs48.exists_point_of_tight}.}

\begin{lemma}[Bing's sequence]\label{lem:bing-seq}
Let $N$ be a nonempty proper subcontinuum. Either there is $p\in X\setminus N$ such that the continua containing $N$ and
missing $p$ have dense union, or some tight proper subcontinuum contains $N$.
\end{lemma}

\begin{proof}
Let $b_0,b_1,\dots$ be a countable basis. Put $N_0=N$, and let $N_{j+1}$ be a proper subcontinuum $K$ such that some
open $O$ has $N_j\cup b_j\subseteq O\subseteq K$, if such a $K$ exists; otherwise $N_{j+1}=N_j$. The $N_j$ form an increasing
sequence of proper subcontinua containing $N$; let $U=\bigcup_jN_j$, a connected set.

\emph{If $U$ is dense, then $U\ne X$.} Suppose $U=X$, and let $W_n$ be the interior of $N_n$; the $W_n$ increase. If
$x\in N_i$, some $j\ge i$ is a step at which $K$ exists (otherwise $N_j=N_i$ for all $j\ge i$ and $N_i=U=X$), and then
$x\in N_j\subseteq O\subseteq N_{j+1}$, so $x\in W_{j+1}$. By compactness $X=W_n$ for some $n$, so $N_n=X$, which is false. Any $p\notin
U$ then satisfies the first alternative, with the continua $N_n$.

\emph{If $U$ is not dense, $N'=\cl U$ is tight.} $N'$ is a proper subcontinuum containing $N$. Let $K$ be a proper
subcontinuum and $O$ open with $N'\subseteq O\subseteq K$. Some $q\in O$ is not in $N'$: otherwise $N'=O$ would be open and closed,
nonempty and not all of $X$. Choose $j$ with $q\in b_j\subseteq O\setminus N'$. Then $N_j\cup b_j\subseteq O\subseteq K$, so at step $j$ a
continuum $N_{j+1}$ was chosen with $b_j\subseteq N_{j+1}$, and $q\in U\subseteq N'$, a contradiction.
\end{proof}

\formal{\lean{HPC.Plane.Arcs48.tight_or_point}, \lean{HPC.Plane.Arcs48.exists_point_dense}.}

\begin{proof}[Proof of Theorem~\ref{thm:bing5}]
$N=\set r$ is a proper subcontinuum. In the first case of Lemma~\ref{lem:bing-seq} we are done; in the second,
apply Lemma~\ref{lem:tight-case} to the tight continuum $N'\ni r$.
\end{proof}

\formal{\lean{HPC.Plane.arcs_theorem5_point}; the general form for a proper subset $R$ in place of $\set r$ is
\lean{HPC.Plane.arcs_theorem5}.}

\begin{corollary}[Bing 1948, Corollary~2]\label{cor:bing2}
A nondegenerate continuum $X$ contains two points $p\ne q$ with the following property: for any nonempty open sets
$D,E\subseteq X$ there is a continuum that meets $D$ and $E$ and misses $p$ (respectively $q$).
\end{corollary}

\begin{proof}
Let $r\in X$ and let $p$ be given by Theorem~\ref{thm:bing5}. Continua $K_1,K_2$ containing $r$ and missing $p$ meet $D$,
respectively $E$; their union is the required continuum. Then let $q\ne p$ be given by Theorem~\ref{thm:bing5} for
$r=p$.
\end{proof}

\formal{\lean{HPC.Plane.arcs_corollary2}, with Bing's relation ``$R$ cuts $A$ from $B$ in $M$'' as
\lean{HPC.Plane.Cuts}.}

\section{Consequences for homogeneous continua}

\begin{proposition}[Jones 1955, claims (3) and (5)]\label{prop:claim35}
Let $X$ be a nondegenerate homogeneous continuum and $p\in X$.
\begin{enumerate}[label=(\alph*)]
\item For any nonempty open sets $D,E$ there is a continuum meeting $D$ and $E$ and missing $p$.
\item Any two points $a,b\in\Us p$ lie in a continuum missing $p$.
\end{enumerate}
\end{proposition}

\begin{proof}
(a) Let $p_0$ be a point given by Corollary~\ref{cor:bing2} and $T$ a homeomorphism with $Tp_0=p$. A continuum
meeting $T^{-1}(D)$ and $T^{-1}(E)$ and missing $p_0$ is carried by $T$ to the required continuum.

(b) Let $A\supseteq V\ni a$ and $B\supseteq W\ni b$ be continua missing $p$ containing open sets $V$, $W$. By (a) a continuum $K$
missing $p$ meets $V$ and $W$. Then $A\cup K\cup B$ is a continuum containing $a$, $b$ and missing $p$.
\end{proof}

\formal{\lean{HPC.Jones.exists_continuum_meets_of_open}, \lean{HPC.Jones.jones55_claim3},
\lean{HPC.Jones.exists_continuum_of_mem_aposSet}, \lean{HPC.Jones.jones55_claim5}.}

\section*{Exercises}

\begin{exercise}
Show that Theorem~\ref{thm:spreading} fails for the arc $[0,1]$ (which is not homogeneous): find a continuum $K$ with
nonempty interior, $o\in K$ and $x\notin K$ such that every continuum containing a neighbourhood of $o$ contains $x$.
\end{exercise}

\begin{exercise}
In the arc $[0,1]$, which points $p$ satisfy the conclusion of Corollary~\ref{cor:bing2}? Show that in the circle every
point does.
\end{exercise}

\begin{exercise}
Prove that a lower semicontinuous function $g\colon X\to[0,\infty]$ on a compact space attains its minimum.
\end{exercise}

%======================================================================
\chapter{Jones's aposyndetic decomposition}\label{ch:jones55}
%======================================================================

This chapter proves Theorem~1 of Jones's paper \emph{On a certain type of homogeneous plane continuum}
(1955) and the part of his Theorem~2 that does not need the 1949 paper. For a homogeneous continuum $X$
the sets $\Ls x$ of Definition~\ref{def:apos} partition $X$; the partition is a continuous decomposition
into terminal continua; and the quotient is a homogeneous aposyndetic continuum. When $X$ lies in the
plane, no element separates the plane.

Jones derives the partition property through twelve numbered claims and a case argument occupying most
of pages 737--738 of his paper. Here it follows in a few lines from the spreading theorem
(Theorem~\ref{thm:spreading}) and the Lemma of Jones 1951 (Theorem~\ref{thm:jones51-lemma}).

Throughout, $X$ is a homogeneous continuum (compact, connected, metric). Parts (b) and (c) of the main
theorem need neither connectedness nor decomposability.

\section{The sets \texorpdfstring{$\Ls x$}{L\_x} partition \texorpdfstring{$X$}{X}}

\begin{lemma}\label{lem:U-eq}
Let $X$ be a homogeneous compact metric space. If $y\in\Ls x$, then $\Us y=\Us x$.
\end{lemma}

\begin{proof}
If $y=x$ there is nothing to prove; otherwise $X$ is not aposyndetic at $y$ with respect to $x$. Let $z\in\Us x$.
If $z\notin\Us y$, then $X$ is not aposyndetic at $z$ with respect to $y$, nor at $y$ with respect to $x$, so by
Corollary~\ref{cor:trans} not at $z$ with respect to $x$, against $z\in\Us x$. Hence $\Us x\subseteq\Us y$. By
Theorem~\ref{thm:jones51-lemma} the inclusion is not proper.
\end{proof}

\formal{\lean{HPC.Jones.aposSet_eq_of_not_aposAt}.}

\begin{theorem}[Jones 1955, Theorem~1(b)]\label{thm:jones55-b}
Let $X$ be a homogeneous compact metric space and $x,y\in X$. Then $y\in\Ls x$ if and only if $\Ls y=\Ls x$.
Consequently any two of the sets $\Ls x$ are equal or disjoint, and they partition $X$.
\end{theorem}

\begin{proof}
If $y\in\Ls x$, then $\Ls y=X\setminus\Us y=X\setminus\Us x=\Ls x$ by Lemma~\ref{lem:U-eq}. Conversely, if
$\Ls y=\Ls x$, then $y\in\Ls y=\Ls x$. If $z\in\Ls x\cap\Ls y$, then $\Ls x=\Ls z=\Ls y$.
\end{proof}

\formal{\lean{HPC.Jones.mem_Lset_iff_Lset_eq}, \lean{HPC.Jones.jones55_b}. Jones assumes $X$ decomposable,
which is not needed.}

\begin{theorem}[Jones 1955, Theorem~1(c): the elements are terminal]\label{thm:jones55-c}
Let $X$ be a homogeneous compact metric space. A continuum $K\subseteq X$ that meets both $\Ls x$ and
$X\setminus\Ls x$ contains $\Ls x$.
\end{theorem}

\begin{proof}
Let $a\in K\cap\Ls x$ and $b\in K\setminus\Ls x$, and suppose $y\in\Ls x\setminus K$. By Theorem~\ref{thm:jones55-b},
$\Ls y=\Ls x$; so $b\notin\Ls y$, that is, some continuum $B$ missing $y$ contains an open set $V\ni b$. Then $K\cup B$
is a continuum (the two continua share $b$) with nonempty interior, it contains $a$ and misses $y$. By
Theorem~\ref{thm:spreading}, $X$ is aposyndetic at $a$ with respect to $y$, so $a\notin\Ls y=\Ls x$, a
contradiction.
\end{proof}

\formal{\lean{HPC.Jones.Lset_subset_of_meets}, \lean{HPC.Jones.jones55_c}.}

So the elements $\Ls x$ are \emph{terminal} subcontinua: a subcontinuum meeting one of them either lies inside
it or contains it. We write $G=\set{\Ls x : x\in X}$ for the decomposition and $\pi\colon X\to X/G$ for the
quotient map; in the notation of \vthree{Chapter~6}, the element of $x$ is $\elt x=\Ls x$.

\formal{The relation $\Ls x=\Ls y$ is \lean{HPC.Jones.lRel}, the setoid is \lean{HPC.Jones.lSetoid}, and
\lean{HPC.Jones.cls_lSetoid} says that its classes are the sets $\Ls x$.}

\begin{remark}[Jones's properties (a)--(f)]\label{rem:props}
On p.~739 Jones lists six properties of an element $g=\Ls o$, phrased with his auxiliary set $N_o$ (the point $o$
together with the points that cut $o$ from $\Us o$). All of them follow from the two theorems above:
$N_o=g$; $\Us x=\Us o$ for $x\in g$ (Lemma~\ref{lem:U-eq}); $N_x=g$ for $x\in g$; two distinct points $x,y$ of $g$ cut
each other from $\Us o$; $X\setminus g=\Us o$ (the definition); and a continuum meeting $g$ and $\Us o$ contains $g$
(Theorem~\ref{thm:jones55-c}). The fourth, for instance: a continuum through $y$ and a point of $\Us o=X\setminus g$
contains $g\ni x$.
\end{remark}

\formal{\lean{HPC.Jones.jones55_prop_a}, \lean{HPC.Jones.jones55_prop_b}, \lean{HPC.Jones.jones55_prop_c},
\lean{HPC.Jones.jones55_prop_d}, \lean{HPC.Jones.jones55_prop_e}, \lean{HPC.Jones.jones55_prop_f}, with $N_o$ as
\lean{HPC.Jones.Nset} and \lean{HPC.Jones.Nset_eq_Lset}. Jones's intermediate claims (1)--(12) are also proved,
as \lean{HPC.Jones.jones55_claim1} through \lean{HPC.Jones.jones55_claim12}.}

\section{The decomposition is continuous}

\begin{theorem}[Jones 1955, p.~739]\label{thm:jones55-usc}
Let $X$ be a homogeneous continuum. The decomposition $G$ is upper semicontinuous, and it is continuous.
\end{theorem}

\begin{proof}
\emph{Upper semicontinuity.} By Theorem~\ref{thm:jones55-b}, $\Ls x=\Ls y$ if and only if $X$ is not aposyndetic
at $x$ with respect to $y$. By Theorem~\ref{thm:jones48-1} the set of such pairs is closed in $X\times X$, and a
decomposition with closed relation is u.s.c.\ (\vthree{Proposition~6.4}).

\emph{Continuity.} By the theorem of Kuratowski and Fort (\vthree{Theorem~6.10}), $G$ is lower semicontinuous at
some element $\Ls{x_0}$. Given $x$, let $T$ be a homeomorphism with $Tx_0=x$. By Lemma~\ref{lem:apos-basic}(b), $T$
maps elements onto elements, $T(\Ls a)=\Ls{Ta}$, so it permutes the elements, and by \vthree{Proposition~6.11} the
decomposition is lower semicontinuous at $\Ls{Tx_0}=\Ls x$.
\end{proof}

\formal{\lean{HPC.Jones.isClosed_lSetoid}, \lean{HPC.Jones.jones55_usc}, \lean{HPC.Jones.lSetoid_homeomorph},
\lean{HPC.Jones.jones55_continuous}; from \lean{HPC.Jones.isUSCDecomp_of_isClosed},
\lean{HPC.Jones.exists_lscAt} (with \lean{HPC.Jones.residual_lscAt}) and \lean{HPC.Jones.lscAt_homeomorph}.
Jones argues upper semicontinuity with convergent sequences of elements and property (f).}

\section{The quotient}

\begin{theorem}[Jones 1955, Theorem~1(a)]\label{thm:jones55-a}
Let $X$ be a decomposable homogeneous continuum. Then $X/G$ is a nondegenerate, homogeneous, aposyndetic,
compact metrizable continuum.
\end{theorem}

\begin{proof}
$X/G$ is compact metrizable by \vthree{Theorem~6.5}, since $G$ is u.s.c., and connected as a continuous image of $X$.

\emph{Nondegenerate.} By Theorem~\ref{thm:jones48-8} there are $x\ne y$ with $X$ aposyndetic at $x$ with respect to $y$.
Then $x\notin\Ls y$, so $\pi(x)\ne\pi(y)$.

\emph{Homogeneous.} Let $a,b\in X$ and $T$ a homeomorphism with $Ta=b$. $T$ permutes the elements, so it induces a
homeomorphism $\hat T$ of $X/G$ with $\hat T(\pi(a))=\pi(b)$ (\vthree{Proposition~6.8}).

\emph{Aposyndetic.} Let $\pi(y)\ne\pi(x)$. Then $y\notin\Ls x$, so some continuum $K$ missing $x$ contains an open $V\ni y$.
By Theorem~\ref{thm:jones55-c}, $K$ misses $\Ls x$: otherwise it meets $\Ls x$ and (at $y$) its complement, so contains
$\Ls x\ni x$. Hence $\pi(K)$ is a continuum in $X/G$ missing $\pi(x)$. By lower semicontinuity at $\Ls y$
(Theorem~\ref{thm:jones55-usc}) there is an open $W\supseteq\Ls y$ such that every element meeting $W$ meets $V$. The
set $W^*$ of points whose element lies in $W$ is open and a union of elements because $G$ is u.s.c.\
(\vthree{Lemma~6.3}), so $\pi(W^*)$ is an open set containing $\pi(y)$; and each element in it meets $V\subseteq K$, so
each point of $\pi(W^*)$ lies in $\pi(K)$. So $\pi(y)\in\pi(W^*)\subseteq\pi(K)$.
\end{proof}

\formal{\lean{HPC.Jones.jones55_a}; homogeneity from \lean{HPC.Jones.exists_homeomorph_quotient},
nondegeneracy from \lean{HPC.Jones.theorem8}. Jones states aposyndesis of $X/G$ without argument; the proof above
uses the continuity of $G$.}

\section{The elements}

\begin{proposition}[Jones 1955, p.~739]\label{prop:jones55-elements}
Let $X$ be a homogeneous continuum.
\begin{enumerate}[label=(\alph*)]
\item If $X$ is aposyndetic, every element is a point: $\Ls x=\set x$.
\item If $X$ is not aposyndetic, every element is nondegenerate, any two elements are homeomorphic, and each element is
a homogeneous continuum.
\item If $X$ is decomposable, every element has empty interior; equivalently, each $\Us x$ is dense.
\end{enumerate}
\end{proposition}

\begin{proof}
(a) is the definition. (b) Some $a\ne b$ have $X$ not aposyndetic at $a$ with respect to $b$, so $a\in\Ls b$ and $\Ls b$
is nondegenerate. A homeomorphism $T$ with $Tb=x$ maps $\Ls b$ onto $\Ls x$; so every element is nondegenerate, and any
two are homeomorphic. Let $a,b\in\Ls x$ and $T$ a homeomorphism with $Ta=b$. By Theorem~\ref{thm:jones55-b},
$T(\Ls x)=T(\Ls a)=\Ls{Ta}=\Ls b=\Ls x$, so $T$ restricts to a homeomorphism of $\Ls x$ carrying $a$ to $b$. Each $\Ls x$
is a continuum by Theorem~\ref{thm:jones48-3}.

(c) $X$ is nondegenerate, and by Theorem~\ref{thm:jones48-8} and homogeneity no $\Us o$ is empty. Suppose the open set
$D=\interior\Ls o$ is nonempty and let $x\in D$. A nonempty open subset of a nondegenerate continuum is not a single
point, so $D\setminus\set x$ is nonempty and open. By Proposition~\ref{prop:claim35}(a) a continuum $K$ missing $x$ meets
$D\setminus\set x\subseteq\Ls o$ and $\Us o=X\setminus\Ls o$. By Theorem~\ref{thm:jones55-c}, $K\supseteq\Ls o\ni x$, a
contradiction.
\end{proof}

\formal{\lean{HPC.Jones.Lset_eq_singleton_of_isAposyndetic}, \lean{HPC.Jones.jones55_elements},
\lean{HPC.Jones.isHomogeneous_Lset}, \lean{HPC.Jones.aposSet_nonempty}, \lean{HPC.Jones.interior_Lset_eq_empty},
\lean{HPC.Jones.jones55_claim12}.}

\begin{proposition}\label{prop:compl-connected}
Let $X$ be a decomposable homogeneous continuum and $x\in X$. Any two points of $X\setminus\Ls x$ lie in a continuum
contained in $X\setminus\Ls x$. In particular $X\setminus\Ls x$ is connected.
\end{proposition}

\begin{proof}
$X\setminus\Ls x=\Us x$. By Proposition~\ref{prop:claim35}(b), two points $a,b\in\Us x$ lie in a continuum $K$ missing $x$.
As $K$ meets $X\setminus\Ls x$ and does not contain $x$, Theorem~\ref{thm:jones55-c} shows that $K$ misses $\Ls x$.
\end{proof}

\formal{\lean{HPC.Jones.jones55_compl_continuumwise}, \lean{HPC.Jones.jones55_compl_connected}; Jones writes
``it is clear that for each $g$ of $G$, $M-g$ is connected (strongly or continuum-wise connected) because of (5)''.}

\section{Elements in the plane do not separate it}

From now on $X=M$ is a homogeneous continuum in the plane $\R^2$.

\begin{lemma}\label{lem:disjoint-open-countable}
A family of pairwise disjoint nonempty open subsets of $\R^2$ is countable.
\end{lemma}

\begin{proof}
Each member contains a point with rational coordinates, and distinct members contain distinct such points.
\end{proof}

\formal{\lean{HPC.Jones.countable_of_pairwiseDisjoint_open}.}

\begin{theorem}[Jones 1955, Theorem~2(2), non-separation]\label{thm:elt-not-separates}
Let $M\subseteq\R^2$ be a decomposable homogeneous continuum. No element $\Ls x$ separates the plane.
\end{theorem}

\begin{proof}
If $M$ is aposyndetic, the elements are points (Proposition~\ref{prop:jones55-elements}(a)), and a point does not
separate the plane. Otherwise the elements are homeomorphic compact sets
(Proposition~\ref{prop:jones55-elements}(b)), so if one separates the plane, all do (\vfour{Theorem~2.16}). Assume this.

For each element $g$ the set $M\setminus g$ is connected (Proposition~\ref{prop:compl-connected}), so it lies in one
component of $\R^2\setminus g$; since $\R^2\setminus g$ has at least two components, some component $W_g$ of $\R^2\setminus g$
misses $M\setminus g$, hence misses $M$. Recall that $\cl{W_g}\subseteq W_g\cup g$ (the closure of a component of the
complement of $g$ adds only points of $g$).

Let $g\ne h$ be elements with $W_g\cap W_h\ne\emptyset$. $W_g$ is a connected subset of $\R^2\setminus h$ (it misses $M\supseteq h$)
meeting the component $W_h$, so $W_g\subseteq W_h$; symmetrically $W_h\subseteq W_g$. Then the closure of $W=W_g=W_h$ lies in
$(W\cup g)\cap(W\cup h)=W$, as $g\cap h=\emptyset$. So $W$ is a nonempty open and closed proper subset of the plane, which is
impossible. Hence the $W_g$ are pairwise disjoint nonempty open sets.

$X/G$ is a nondegenerate metrizable continuum (Theorem~\ref{thm:jones55-a}), hence uncountable: for two of its points
$p\ne q$, the continuous function $\dist(p,\cdot)$ takes every value between $0$ and $\dist(p,q)$. So there are
uncountably many sets $W_g$, contradicting Lemma~\ref{lem:disjoint-open-countable}.
\end{proof}

\formal{\lean{HPC.Jones.jones55_elt_not_separates}, using \lean{HPC.Jones.separatesPlane_of_homeomorph},
\lean{HPC.Jones.closure_component_compl_subset}, \lean{HPC.Jones.uncountable_quotient} and
\lean{HPC.Jones.uncountable_of_connected_metrizable}. Jones's proof says only that ``if one separates, all do, which
is impossible''; the disjointness of the domains $W_g$ is the missing step.}

\section{An application: Jones's 1951 theorem}

\begin{theorem}[Jones 1951, Theorem~1]\label{thm:jones51-1}
A homogeneous hereditarily unicoherent continuum is indecomposable.
\end{theorem}

\begin{proof}
Suppose $X$ is decomposable. By Theorem~\ref{thm:jones48-8} and homogeneity there are $a,b$ with $b\in\Us a$, so
$\Ls a\neq\Ls b$, and the two elements are disjoint. Let $J$ be a subcontinuum irreducible between $a$ and $b$
(\vthree{Theorem~2.5}). By hereditary unicoherence, $J$ lies in every continuum $K$ containing $a$ and $b$: $J\cap K$ is a
continuum containing $a$ and $b$, so equals $J$. The connected set $J$ is not covered by the two disjoint closed sets
$\Ls a$, $\Ls b$, so some $c\in J$ lies in neither. Then $a,b\in\Us c$, and by Proposition~\ref{prop:claim35}(b) some
continuum containing $a$ and $b$ misses $c$. It contains $J\ni c$, a contradiction.
\end{proof}

\formal{\lean{HPC.Jones.jones51_theorem1}. Jones's own proof builds the decomposition of the 1955 paper in an
early form, passes to the decomposition space and applies Whyburn's (6.21) there.}

\section*{Exercises}

\begin{exercise}
Show that the conclusion of Theorem~\ref{thm:jones55-b} fails for the harmonic fan (it is not homogeneous).
\end{exercise}

\begin{exercise}
Let $X=P\times\Sph$, where $P$ is a pseudo-arc. Assuming that $X$ is homogeneous, determine the sets $\Ls x$ and the
quotient $X/G$. (Compare Theorem~\ref{thm:jones48-7}: $X$ is aposyndetic.) Do the same for the circle of pseudo-arcs
of Chapter~\ref{ch:decomposable}, using Theorem~\ref{thm:jones55-c}.
\end{exercise}

\begin{exercise}
Show that every subcontinuum of a hereditarily indecomposable continuum is terminal. Let $X$ be the union of a
nondegenerate indecomposable continuum $K$ and an arc $A$ with $A\cap K$ an end point of $A$. Show that $\Ls y\supseteq K$ for
$y\in K$, so that Proposition~\ref{prop:jones55-elements}(c) fails without homogeneity.
\end{exercise}

\begin{exercise}
Where does the proof of Theorem~\ref{thm:elt-not-separates} use that $M$ is decomposable?
\end{exercise}

%======================================================================
\chapter{Order and Menger's theorem}\label{ch:order}
%======================================================================

The second half of the volume proves Jones's theorem of 1949: a nondegenerate homogeneous aposyndetic plane
continuum is a simple closed curve. Its proof begins: ``If a point of $M$ is of order~2 in $M$, then $M$ is a
simple closed curve.'' This chapter supplies that sentence. It introduces the Menger--Urysohn order of a point
and proves the case of Menger's theorem that is needed: a nondegenerate continuum all of whose points have
order~2 is a simple closed curve.

\section{Order}

\begin{definition}\label{def:order}
Let $Y$ be a topological space, $p\in Y$ and $n\in\N$. The point $p$ has \emph{order at most $n$} if every
neighbourhood $U$ of $p$ contains an open set $V\ni p$ whose frontier $\fr V$ is finite with at most $n$ points. It has
\emph{order $n$} if it has order at most $n$ but not order at most $n-1$.
\end{definition}

\formal{\lean{HPC.Jones.OrderLE}, \lean{HPC.Jones.HasOrder}.}

For example, every point of a circle has order~2, the end points of an arc have order~1, and the vertex of a
simple triod has order~3. In a nondegenerate continuum no point has order~$0$: a small open $V\ni p$ with empty
frontier would be open and closed.

\begin{lemma}\label{lem:order-basic}
\leavevmode
\begin{enumerate}[label=(\alph*)]
\item Order is a topological invariant: if $h\colon Y\to Z$ is a homeomorphism, $p$ has order at most $n$ (order $n$)
in $Y$ if and only if $h(p)$ does in $Z$.
\item If $p\in W\subseteq Y$ and $p$ has order at most $n$ in $Y$, then it has order at most $n$ in the subspace $W$.
\end{enumerate}
\end{lemma}

\begin{proof}
(a) $h$ carries open sets, neighbourhoods and frontiers to the same. (b) Given a neighbourhood $U\cap W$ of $p$ in $W$
with $U$ open in $Y$, take $V\subseteq U$ for $Y$; the frontier of $V\cap W$ in $W$ lies in $\fr V$.
\end{proof}

\formal{\lean{HPC.Jones.orderLE_homeomorph}, \lean{HPC.Jones.hasOrder_homeomorph}, \lean{HPC.Jones.orderLE_subtype}.}

\begin{lemma}[Triods raise the order]\label{lem:triod}
Let $Y$ be a $T_1$ space and $T_1,T_2,T_3$ connected sets containing $p$, each containing a point other than $p$,
with $T_i\cap T_j\subseteq\set p$ for $i\ne j$. Then $p$ does not have order at most~$2$.
\end{lemma}

\begin{proof}
Choose $x_i\in T_i\setminus\set p$ and let $V\ni p$ be open, missing the finite closed set $\set{x_1,x_2,x_3}$. Each $T_i$
is connected and contains $p\in V$ and $x_i\notin V$, so it meets $\fr V$, at a point $f_i\ne p$ (as $p\in V$ and $V\cap\fr
V=\emptyset$). The $f_i$ are distinct because the $T_i$ meet only at $p$. So $\fr V$ has at least three points.
\end{proof}

\formal{\lean{HPC.Jones.not_orderLE_two_of_triod}.}

\section{Curves of finite order}

\begin{proposition}\label{prop:finite-order-lc}
A continuum $Y$ each of whose points has finite order is locally connected.
\end{proposition}

\begin{proof}
Let $p\in Y$ and $\varepsilon>0$. Choose $n$ with $p$ of order at most $n$, and an open $V\ni p$ inside $\bar B(p,\varepsilon/2)$
with finite frontier. Then $F=\cl V\subseteq B(p,\varepsilon)$. If $F=Y$, the connected set $Y$ is a neighbourhood of $p$ inside
$B(p,\varepsilon)$. Otherwise, by boundary bumping (\vthree{Theorem~1.10}), the component in $F$ of each $z\in F$ meets
$\cl{Y\setminus F}$ at a point $w$; and $w\notin V$ because $V$ is open and $V\subseteq\interior F$; so $w\in\fr V$. Hence the
components of $F$ through the finitely many points of $\fr V$ cover $F$. The union $D$ of those not containing $p$ is
closed, and $V\setminus D$ is an open set containing $p$ and contained in the component $C$ of $p$ in $F$. So the connected set $C$
is a neighbourhood of $p$ inside $B(p,\varepsilon)$.
\end{proof}

\formal{\lean{HPC.Jones.locallyConnectedSpace_of_orderLE}.}

\begin{lemma}\label{lem:open-continuum}
In a compact, locally connected Hausdorff space, any two points of a connected open set $O$ lie in a continuum
contained in $O$.
\end{lemma}

\begin{proof}
Fix $x\in O$ and let $S$ be the set of $z\in O$ such that some continuum $K\subseteq O$ contains $x$ and $z$. Every $z\in O$ has a
connected open neighbourhood $N$ with $\cl N\subseteq O$; $\cl N$ is a continuum, so either all points of $N$ are in $S$ (use
$K\cup\cl N$) or none is. Thus $S$ and $O\setminus S$ are open; $x\in S$, and $O$ is connected, so $S=O$.
\end{proof}

\formal{\lean{HPC.Jones.exists_continuum_of_isOpen}.}

\begin{proposition}[Arcs in curves of finite order]\label{prop:finite-order-arcs}
Let $Y$ be a compact metric space each of whose points has finite order. Any two distinct points $a,b$ of a
subcontinuum $K$ of $Y$ are the end points of an arc in $K$.
\end{proposition}

\begin{proof}
Let $J\subseteq K$ be a subcontinuum irreducible between $a$ and $b$ (\vthree{Theorem~2.5}). By Lemma~\ref{lem:order-basic}(b) each
point of $J$ has finite order in $J$, so $J$ is locally connected (Proposition~\ref{prop:finite-order-lc}). Let $c\in
J\setminus\set{a,b}$. If $J\setminus\set c$ were connected, then by Lemma~\ref{lem:open-continuum} (applied in $J$) a continuum
inside $J\setminus\set c$ would contain $a$ and $b$, against irreducibility. So every point of $J$ other than $a,b$ cuts $J$: $J$ is a
cut-point arc from $a$ to $b$, hence an arc (\vfive{Proposition~1.11}).
\end{proof}

\formal{\lean{HPC.Jones.exists_tarc_of_orderLE}.}

\begin{lemma}[An arc with legs at both ends]\label{lem:legs}
Let $Y$ be a compact metric space all of whose points have order at most $2$, $K\subseteq Y$ a continuum and $B\subseteq K$ an arc
from $a$ to $b$. Suppose there are connected sets $T_a\ni a$ and $T_b\ni b$, each with a second point, such that $T_a\cap
K\subseteq\set a$ and $T_b\cap K\subseteq\set b$. Then $K=B$.
\end{lemma}

\begin{proof}
Suppose $P\in K\setminus B$. By Proposition~\ref{prop:finite-order-arcs} there is an arc in $K$ from $P$ to $a$; let $q$ be its first point
on $B$, and $C$ the sub-arc from $P$ to $q$, so $C\cap B=\set q$ (\vfive{Lemma~1.4}). If $q=a$, the connected sets $C$, $B$, $T_a$ form a
triod at $a$: they meet pairwise only in $a$ because $C,B\subseteq K$. This contradicts Lemma~\ref{lem:triod}. The case $q=b$ is the same with
$T_b$. If $q$ is an interior point of $B$, the two sub-arcs of $B$ from $a$ to $q$ and from $q$ to $b$ together with $C$ form a triod at $q$.
\end{proof}

\formal{\lean{HPC.Jones.arc_eq_of_legs}.}

\section{Menger's theorem}

\begin{theorem}[Menger]\label{thm:menger}
A nondegenerate continuum $X$ all of whose points have order $2$ is homeomorphic to the circle.
\end{theorem}

\begin{proof}
$X$ is locally connected by Proposition~\ref{prop:finite-order-lc}. Let $u$ be a non-cut point of $X$ (\vthree{Theorem~3.6}). As $u$ is
not of order at most $1$, some neighbourhood $U_0$ of $u$ contains no open $V\ni u$ with at most one frontier point. Let $V_0\ni u$ be
open, with $\cl{V_0}\subseteq U_0$, inside a ball about $u$ of radius less than half the diameter of $X$, with at most two frontier points, and
let $V$ be the component of $u$ in $V_0$. In a locally connected space $V$ is open and $\fr V\subseteq\fr V_0$. By the choice of $U_0$,
$\fr V=\set{a,b}$ with $a\ne b$; and $\cl V\ne X$ by the choice of radius.

\emph{$\cl V$ is an arc.} $\cl V$ is a continuum, so it contains an arc $B$ from $a$ to $b$ (Proposition~\ref{prop:finite-order-arcs}). Take
connected open neighbourhoods $N_a$ of $a$ and $N_b$ of $b$ missing $b$, respectively $a$. Since $N_a$ meets $\fr V$ only in $a$, the set
$T_a=N_a\setminus V$ is connected, and it contains points outside $\cl V$ (as $a\in\cl{X\setminus\cl V}$, for otherwise $a$ would be
interior to $\cl V$, and $V\cup\set a\subseteq U_0$ would be open with frontier $\set b$). Likewise for $T_b$. Lemma~\ref{lem:legs} gives
$\cl V=B$. As $u\in V$ is not an end point, $B\setminus\set u$ consists of two half-arcs $[a,u)$ and $(u,b]$.

\emph{$X\setminus V$ is connected.} Suppose $X\setminus V=P\cup Q$ with $P,Q$ closed, disjoint and nonempty. If $a,b\in P$, then $Q$ is closed and
also open in $X\setminus\set u$, because its complement there is $P\cup(V\setminus\set u)$, whose closure lies in $P\cup V$. If $a\in P$ and
$b\in Q$, then $X\setminus\set u=(P\cup[a,u))\cup(Q\cup(u,b])$ is a separation. Both contradict that $u$ is not a cut point.

\emph{$X\setminus V$ is an arc.} It is a continuum containing $a$ and $b$; let $B'\subseteq X\setminus V$ be an arc from $a$ to $b$. The legs
$[a,u]$ and $[u,b]$ (the two halves of $\cl V$) meet $X\setminus V$ only in $a$, respectively $b$. By Lemma~\ref{lem:legs}, $X\setminus V=B'$.

So $X=B\cup B'$ is the union of two arcs from $a$ to $b$ with $B\cap B'=\cl V\cap(X\setminus V)=\set{a,b}$: a simple closed curve,
homeomorphic to the circle (\vfive{Proposition~1.13}).
\end{proof}

\formal{\lean{HPC.Jones.menger_order_two}, with \lean{HPC.Jones.frontier_connectedComponentIn_subset},
\lean{HPC.Jones.isPreconnected_diff_of_frontier} and \lean{HPC.Jones.tscc_homeomorph_circle} (two arcs glued along their ends,
\lean{HPC.Jones.glue_circle}). Menger's own route goes through his $n$-Beinsatz and the classification of curves of order~2; the
proof above is his second argument (the first point on an arc is a branch point), applied twice.}

\begin{corollary}[Jones 1949, first sentence of the proof]\label{cor:order-two}
If a homogeneous continuum $X$ has a point of order $2$, it is homeomorphic to the circle.
\end{corollary}

\begin{proof}
By homogeneity and Lemma~\ref{lem:order-basic}(a), every point has order $2$. $X$ is nondegenerate, since a point has order $0$ in a
one-point space. Apply Theorem~\ref{thm:menger}.
\end{proof}

\formal{\lean{HPC.Jones.jones49_order_two}.}

\section*{Exercises}

\begin{exercise}
Show that every point of the $\sin(1/x)$-continuum lying on the limit segment has no finite order, while all other points have
order at most $2$.
\end{exercise}

\begin{exercise}
Prove the converse of Theorem~\ref{thm:menger}: every point of a simple closed curve has order $2$.
\end{exercise}

\begin{exercise}
Show that a continuum all of whose points have order at most $2$ is an arc or a simple closed curve. (Use
Lemma~\ref{lem:legs} with an arc between two points of order $1$.)
\end{exercise}

%======================================================================
\chapter{Whyburn's separating point--order theorem}\label{ch:whyburn}
%======================================================================

Jones's 1949 proof has two cases, and both end the same way: some point of $M$ turns out to be a local separating
point, and then Whyburn's theorem (\emph{Analytic Topology}, Chapter~III, (9.2)) produces a point of order~$2$, to
which Corollary~\ref{cor:order-two} applies. Whyburn's theorem says that in a separable, locally compact metric space
all but countably many local separating points have order~$2$. Whyburn derives it from his theory of cyclic
elements and of ``potential order'' relative to a collection of sets (Chapter~III, \S\S1--2, (8.1)--(8.4)). The
formal proof extracts the part that is needed: a chain of separating points, ordered as Whyburn orders them, is
countable unless its points have order at most~$2$.

Throughout, $X$ is a metric space. From Section~\ref{sec:chains} on it is separable and locally compact.

\section{Separating points of an open set}

\begin{definition}\label{def:sep-point}
Let $R\subseteq X$ and $p\in R$, and let $Q$ be the component of $R$ containing $p$.
\begin{enumerate}[label=(\alph*)]
\item $p$ is a \emph{separating point of $R$} if there are open sets $A$, $B$ with $R\setminus\set p\subseteq A\cup B$, $A$ and $B$
disjoint on $R\setminus\set p$, and both $A$ and $B$ meeting $Q\setminus\set p$.
\item $p$ is a \emph{local separating point} if it is a separating point of some open set $R$.
\end{enumerate}
\end{definition}

\formal{\lean{HPC.Jones.IsSepPointOf}, \lean{HPC.Jones.IsLocalSepPoint}. Whyburn's phrase is ``$R-p=R_a+R_b$ with
$a\in R_a$, $b\in R_b$ in the component of $R$ containing $p$''.}

A cut point $p$ of a connected space $X$ is a separating point of $R=X$. Every point of a circle, and every interior point
of an arc, is a local separating point. The definition is invariant under homeomorphisms.

\formal{\lean{HPC.Jones.IsLocalSepPoint.homeomorph}.}

\begin{lemma}\label{lem:side}
Let $Q$ be connected, $p\in Q$, and let $A$, $B$ be open sets covering $Q\setminus\set p$ and disjoint on it. Then
$Q\setminus(B\setminus\set p)$, the $A$-side of $Q$ together with $p$, is connected.
\end{lemma}

\begin{proof}
Let $D=Q\setminus(B\setminus\set p)$; then $D\setminus\set p\subseteq A$. The set $Q\cap B\setminus\set p$ is open in $Q$, so $D$ is closed in $Q$.
Suppose $D=D_1\cup D_2$ with $D_1$, $D_2$ disjoint, closed in $D$ and $p\in D_1$. Then $D_2$ is closed in $Q$. It is also open in $Q$: if $O$ is
open with $D_2=D\cap O$, we may assume $p\notin O$, and then $D_2=Q\cap(O\cap A)$, because every point of $Q\cap A$ other than $p$ lies in $D$. As $Q$ is
connected and $p\notin D_2$, $D_2=\emptyset$.
\end{proof}

\formal{\lean{HPC.Jones.isPreconnected_side}, \lean{HPC.Jones.isPreconnected_diff_side}.}

\begin{proposition}[Whyburn (9.1)]\label{prop:whyburn91}
Let $X$ be locally compact. If $p$ is a separating point of $R$, it is a separating point of every open $U$ with $p\in U\subseteq R$.
\end{proposition}

\begin{proof}
Let $A$, $B$ be the open sets for $R$ and $Q$ the component of $p$ in $R$. The same $A$, $B$ cover $U\setminus\set p$ and are disjoint on it;
it remains to see that the component $Q_U$ of $p$ in $U$ meets $A\setminus\set p$ (and likewise $B\setminus\set p$). The set $D=Q\setminus(B\setminus\set p)$ is
connected (Lemma~\ref{lem:side}), contains $p$ and a point of $A\setminus\set p$. Let $K\subseteq U$ be a compact neighbourhood of $p$. If $D\subseteq K$, then $D$ is
a connected subset of $U$ through $p$, so $D\subseteq Q_U$ and we are done. Otherwise $D\cap K$ is a compact proper subset of $D$ (it is closed in $D$,
as $D$ is closed in $R$), and by boundary bumping some continuum $C\subseteq D\cap K$ through $p$ meets $\cl{D\setminus K}$, so contains a point other than
$p$ (as $K$ is a neighbourhood of $p$). Then $C\subseteq Q_U$, and $C\setminus\set p\subseteq D\setminus\set p\subseteq A$.
\end{proof}

\formal{\lean{HPC.Jones.whyburn_9_1}, with \lean{HPC.Jones.exists_continuum_bump_compact} (boundary bumping for a compact
piece of a connected set) and \lean{HPC.Jones.whyburn_9_1_side}. Whyburn uses at this point that ``the components
\dots\ are nondegenerate'' (his I, (10.2)).}

\begin{lemma}\label{lem:sep-not-order1}
A separating point $p$ of a set $R$ does not have order at most $1$. Hence, if it has order at most $2$, it has order $2$.
\end{lemma}

\begin{proof}
Let $a\in A\cap Q$ and $b\in B\cap Q$ be points other than $p$, and let $V\ni p$ be open and miss $a$ and $b$. The connected set
$Q\setminus(B\setminus\set p)$ contains $p\in V$ and $a\notin V$, so it meets $\fr V$, at a point of $A$ other than $p$. Likewise the $B$-side gives a
frontier point in $B$. These two points are distinct, so $\fr V$ has at least two points.
\end{proof}

\formal{\lean{HPC.Jones.not_orderLE_one_of_isSepPointOf}, \lean{HPC.Jones.hasOrder_two_of_isSepPointOf}.}

\section{Separation chains}\label{sec:chains}

\begin{definition}\label{def:sepchain}
A \emph{separation chain} consists of a set $R\subseteq X$, a connected set $Q\subseteq R$ that is closed in $R$, a set $T\subseteq Q$, and for each
$p\in T$ two open sets $A_p$, $B_p$ with $A_p\cup B_p=R\setminus\set p$ and $A_p\cap B_p=\emptyset$, such that any two of the sets $A_p$ (and any two of
the sets $B_p$) have a common point in $Q$. For $p,q\in T$ write $p<q$ if $q\in B_p$.
\end{definition}

\formal{\lean{HPC.Jones.SepChain}, with the order \lean{HPC.Jones.SepChain.lt}.}

This is the situation of Whyburn's (1.3): all points of $T$ separate the same two sets in $R$.

\begin{lemma}[Whyburn (1.2), (1.31)]\label{lem:chain-order}
\leavevmode
\begin{enumerate}[label=(\alph*)]
\item For distinct $p,q\in T$: $p<q$ if and only if $p\in A_q$. The relation $<$ is a strict linear order on $T$.
\item If $p<q$, then $Q\cap B_q\subseteq B_p$ and $Q\cap A_p\subseteq A_q$.
\end{enumerate}
\end{lemma}

\begin{proof}
By Lemma~\ref{lem:side}, for $p\in T$ the sets $Q\setminus B_p$ and $Q\setminus A_p$ (the two sides of $p$ in $Q$, each with $p$) are connected. A connected
subset of $R\setminus\set q$ meeting $B_q$ lies in $B_q$, and similarly for $A_q$.

(a) Let $q\in B_p$ and suppose $p\in B_q$. The connected set $Q\setminus B_p$ misses $q$ and contains $p\in B_q$, so it lies in $B_q$. But it contains $Q\cap
A_p$, which meets $Q\cap A_q$; this point would be in $A_q\cap B_q$. So $p\in A_q$. The converse is the same argument with the roles of $A$ and $B$
exchanged. Totality: $q\in R\setminus\set p=A_p\cup B_p$, and $q\in A_p$ means $q<p$. Asymmetry: $q\in B_p$ and $p\in B_q$ is excluded by (a). Transitivity
follows from (b): if $p<q<r$, then $r\in Q\cap B_q\subseteq B_p$.

(b) The connected set $Q\setminus A_q$ misses $p\in A_q$ and contains $q\in B_p$, so it lies in $B_p$; it contains $Q\cap B_q$. The second inclusion is
symmetric.
\end{proof}

\formal{\lean{HPC.Jones.SepChain.lt_iff}, \lean{HPC.Jones.SepChain.lt_total}, \lean{HPC.Jones.SepChain.lt_asymm},
\lean{HPC.Jones.SepChain.lt_trans}, \lean{HPC.Jones.SepChain.B_sub}, \lean{HPC.Jones.SepChain.A_sub}.}

For $x<y$ in $T$ let $M(x,y)=Q\cap\bigl((B_x\cap A_y)\cup\set{x,y}\bigr)$, the \emph{interval} from $x$ to $y$. It is connected and closed in $R$;
intervals shrink as $x$ increases or $y$ decreases; and if $x<x'<y'<y$ then $M(x',y')\subseteq B_x\cap A_y$.

\formal{\lean{HPC.Jones.SepChain.M}, \lean{HPC.Jones.SepChain.M_isPreconnected}, \lean{HPC.Jones.SepChain.M_closed},
\lean{HPC.Jones.SepChain.M_mono_left}, \lean{HPC.Jones.SepChain.M_mono_right}, \lean{HPC.Jones.SepChain.M_sub_BA}.}

\begin{lemma}[The counting argument of Whyburn (2.1)]\label{lem:counting}
Let $X$ be second countable, $T\subseteq X$, $<$ a relation comparing any two distinct points of $T$, $g\colon T\to X$ and $P\colon T\to 2^X$ with
$g(p)\in P(q)$ whenever $p<q$. Then only countably many $p\in T$ have $g(p)\notin\cl{P(p)}$.
\end{lemma}

\begin{proof}
Assign to each such $p$ a basic open set $b(p)\ni g(p)$ missing $P(p)$. If $b(p)=b(q)$ with $p\ne q$, say $p<q$, then $g(p)\in P(q)$, but $g(p)\in b(p)=b(q)$,
which misses $P(q)$. So $b$ is injective, and there are countably many basic sets.
\end{proof}

\formal{\lean{HPC.Jones.countable_not_mem_closure_of_lt}.}

Call $p\in T$ \emph{good} if it is a limit point of its predecessors $\set{q\in T:q<p}$ and of its successors $\set{q\in T:p<q}$. The
\emph{residue} $\operatorname{Res}(p)$ of $p$ is the set of points $z\in Q\setminus\set p$ that lie in $B_x\cap A_y$ for all $x<p<y$ in $T$: the points not
separated from $p$ by the chain.

\formal{\lean{HPC.Jones.SepChain.Good}, \lean{HPC.Jones.SepChain.Res}.}

\begin{lemma}\label{lem:good-res}
Let $X$ be second countable.
\begin{enumerate}[label=(\alph*)]
\item All but countably many points of $T$ are good.
\item All but countably many good points have empty residue.
\end{enumerate}
\end{lemma}

\begin{proof}
(a) Apply Lemma~\ref{lem:counting} with $g(p)=p$ and $P(p)$ the set of predecessors of $p$; then with the reversed order and the successors.

(b) If $p<q$ and $q$ is good, then $\operatorname{Res}(q)\subseteq B_p$: there is $x\in T$ with $p<x<q$ (as $q$ is a limit of predecessors), so $\operatorname{Res}(q)\subseteq
Q\cap B_x\subseteq B_p$ by Lemma~\ref{lem:chain-order}(b). Similarly $\operatorname{Res}(q)\subseteq A_p$ if $q<p$. For each good $p$ with nonempty residue choose
$z(p)\in\operatorname{Res}(p)$; it lies in $A_p$ or in $B_p$. Those $p$ with $z(p)\in B_p$ satisfy $z(p)\notin\cl{A_p}$ ($B_p$ is an open set missing $A_p$), and
$z(p)\in A_q$ whenever $p<q$; Lemma~\ref{lem:counting} with $P(q)=A_q$ shows that they are countably many. The others are handled with the reversed order and
$P(q)=B_q$.
\end{proof}

\formal{\lean{HPC.Jones.SepChain.countable_not_good}, \lean{HPC.Jones.SepChain.Res_sub_B},
\lean{HPC.Jones.SepChain.Res_sub_A}, \lean{HPC.Jones.SepChain.countable_res}.}

\begin{proposition}\label{prop:good-order2}
Let $X$ be locally compact and $p\in T$ a good point with empty residue, interior to $R$. Then $p$ has order at most $2$.
\end{proposition}

\begin{proof}
Let $\rho>0$ be so small that the closed ball $\bar B(p,\rho)$ is compact and lies in $R$. We find $x<p<y$ and an open $V\ni p$ inside $B(p,\rho)$ with
$\fr V\subseteq\set{x,y}$.

\emph{Small intervals.} Some interval $M(x,y)$ with $x<p<y$ lies in $B(p,\rho)$. Otherwise each $M(x,y)$, a connected set through $p$ reaching outside $B(p,\rho)$,
meets the sphere $S=\set{z:d(p,z)=\rho/2}$. The sets $M(x,y)\cap S$ are compact, nonempty and directed downwards (for two pairs, take the larger $x$ and the
smaller $y$), so they have a common point $z\ne p$. For $x<p<y$ choose $x'$, $y'$ in $T$ with $x<x'<p<y'<y$ ($p$ is good); then $z\in M(x',y')\subseteq B_x\cap A_y$.
So $z\in\operatorname{Res}(p)=\emptyset$, a contradiction.

\emph{Cut wire.} Let $O=B_x\cap A_y$, an open set containing $p$, and $F=\bar B(p,\rho)\cap\cl O$, a compact subset of $R$. A point of $\cl O\cap R$ not in $O$
is $x$ or $y$. The component of $p$ in $F$ is a connected subset of $R$ through $p$, so lies in $Q$, hence in $M(x,y)\subseteq B(p,\rho)$; it misses the
closed set $H=\set{z:d(p,z)=\rho}$. By the Shura-Bura cut-wire lemma there is an open $G\ni p$ with $F\cap G$ closed and disjoint from $H$. Put $V=G\cap
O\cap B(p,\rho)$. Its closure lies in the closed set $F\cap G$, which misses $H$, so lies in $B(p,\rho)$; a point of $\cl V\cap O$ is in $G$, hence in $V$.
So $\fr V\subseteq\set{x,y}$.
\end{proof}

\formal{\lean{HPC.Jones.SepChain.exists_small_interval}, \lean{HPC.Jones.SepChain.exists_open_frontier_pair} (with the cut-wire
lemma \lean{HPC.Jones.exists_clopen_of_connectedComponentIn_disjoint}), \lean{HPC.Jones.SepChain.orderLE_two}. This replaces
Whyburn's (8.3) and (2.21)--(2.22).}

\begin{corollary}\label{cor:chain-countable}
In a separable, locally compact metric space, a separation chain with $R$ open none of whose points has order at most $2$ is countable.
\end{corollary}

\begin{proof}
By Proposition~\ref{prop:good-order2} each point of $T$ is not good or is good with nonempty residue; apply Lemma~\ref{lem:good-res}.
\end{proof}

\formal{\lean{HPC.Jones.SepChain.countable_of_not_orderLE_two}.}

\section{The theorem}

\begin{theorem}[Whyburn's separating point--order theorem]\label{thm:sep-order}
Let $X$ be a separable, locally compact metric space and $R\subseteq X$ open. All but countably many separating points of $R$ have order $2$.
\end{theorem}

\begin{proof}
By Lemma~\ref{lem:sep-not-order1} it suffices to count the separating points that do not have order at most $2$. For such a point $p$, take $A$, $B$ as in
Definition~\ref{def:sep-point} (cut down to $R\setminus\set p$) and points $a\in A$, $b\in B$ of the component $Q_p$ of $p$ in $R$, other than $p$; choose basic open sets
$U\ni a$, $W\ni b$ with $U\subseteq A$, $W\subseteq B$. There are countably many pairs $(U,W)$. Fix one, and let $T$ be the set of points $p$ obtained with it.

\emph{$T$ lies in one component of $R$.} Let $p,q\in T$ be distinct. The component $Q_q$ meets $U\subseteq A_p$ and $W\subseteq B_p$. If $p\notin Q_q$, the connected set $Q_q$
lies in $R\setminus\set p=A_p\cup B_p$ and meets both, which is impossible. So $p\in Q_q$.

Hence, with $Q$ this component (connected, and closed in $R$), $T$ and the sets $A_p$, $B_p$ form a separation chain: any two sets $A_p$ contain the nonempty set
$U\cap Q$, and any two sets $B_p$ contain $W\cap Q$. By Corollary~\ref{cor:chain-countable}, $T$ is countable.
\end{proof}

\formal{\lean{HPC.Jones.countable_sepPoint_not_hasOrder_two}, from \lean{HPC.Jones.countable_sepPoint_not_orderLE_two},
\lean{HPC.Jones.sepT}, \lean{HPC.Jones.sepT_mem_component}, \lean{HPC.Jones.countable_sepT}.}

\begin{corollary}[Whyburn (3.2), (8.4), (9.2)]\label{cor:whyburn}
\leavevmode
\begin{enumerate}[label=(\alph*)]
\item In a separable, locally compact metric space only countably many local separating points fail to have order $2$.
\item In a compact metric space only countably many separating points of the whole space fail to have order $2$.
\item In a continuum only countably many cut points fail to have order $2$.
\end{enumerate}
\end{corollary}

\begin{proof}
(a) A separating point $p$ of an open $R$ is, by Proposition~\ref{prop:whyburn91}, a separating point of a basic open set $U$ with $p\in U\subseteq R$. Apply
Theorem~\ref{thm:sep-order} to each of the countably many basic sets. (b) is Theorem~\ref{thm:sep-order} with $R=X$. (c) A cut point $p$ of a continuum $X$ is a
separating point of $X$, unless $X\setminus\set p$ is empty, which happens for at most one point.
\end{proof}

\formal{\lean{HPC.Jones.whyburn_9_2}, \lean{HPC.Jones.whyburn_8_4}, \lean{HPC.Jones.whyburn_3_2}.}

\begin{corollary}[The consequence Jones uses]\label{cor:some-order-two}
An uncountable homogeneous continuum with a local separating point has a point of order $2$.
\end{corollary}

\begin{proof}
By homogeneity every point is a local separating point. If none had order $2$, Corollary~\ref{cor:whyburn}(a) would make the space countable.
\end{proof}

\formal{\lean{HPC.Jones.jones49_some_order_two}, \lean{HPC.Jones.jones49_some_order_two_of_whyburn_9_2}.}

\section*{Exercises}

\begin{exercise}
Show that the Sierpi\'nski triangle (gasket) has exactly three points of order $2$, its corners. Conclude from Corollary~\ref{cor:whyburn} that only countably many of its
points are local separating points, and find infinitely many that are.
\end{exercise}

\begin{exercise}
Show that the vertex of a simple triod is a separating point of the triod that does not have order $2$. Construct a continuum with infinitely many cut
points, none of which has order $2$.
\end{exercise}

\begin{exercise}
In the arc $[0,1]$ with $R=(0,1)$, $Q=R$, $T=\set{1/n:n\ge2}$, $A_p=(0,p)$ and $B_p=(p,1)$, check that the data form a separation chain, and show
that none of its points is good. Find a chain in $[0,1]$ all of whose points are good.
\end{exercise}

%======================================================================
\chapter{Continua crossing an annulus}\label{ch:sectors}
%======================================================================

Whyburn's proof that the boundary of a complementary domain of a semi-locally connected plane continuum is
locally connected (Chapter~\ref{ch:boundaries}) argues by contradiction: if the boundary is not locally
connected, infinitely many disjoint pieces of it cross an annulus. Whyburn, following Moore (1922), then takes
for granted two facts about such crossing continua: two of them cut the annulus into two pieces, and
infinitely many of them can be arranged in a cyclic order, each one cut off from the others by its two
neighbours. This chapter proves both facts, using only continuous logarithms (\vfour{Chapters~1--2}): neither
the Jordan curve theorem nor Janiszewski's theorem is needed.

The idea is to unroll the annulus. Polar coordinates map a strip onto the annulus; a crossing continuum that
misses another crossing continuum lifts to a crossing set of the strip; in the strip, crossing sets are
linearly ordered from left to right; and the order on the lifts gives the cyclic order.

Throughout, the plane $\R^2$ carries the maximum metric $d(z,w)=\max(\abs{z_1-w_1},\abs{z_2-w_2})$, so that ``circles''
are squares.

\section{Crossing continua and the strip}

\begin{definition}\label{def:annulus}
Let $P\in\R^2$ and $0<r_1<r_2$.
\begin{enumerate}[label=(\alph*)]
\item The \emph{square annulus} is $\Ann=\Ann(P,r_1,r_2)=\set{z: r_1\le d(z,P)\le r_2}$. A continuum $K$ \emph{crosses} it if
$K\subseteq\Ann$ and $K$ meets both squares $d(z,P)=r_1$ and $d(z,P)=r_2$.
\item The \emph{strip} is $\Strip=\Strip(a,b)=\set{(x,y): a\le y\le b}$. A \emph{crossing set} of the strip is a continuum in $\Strip$
meeting both edges $y=a$ and $y=b$.
\item The \emph{covering map} $E\colon\Strip(r_1,r_2)\to\Ann$ is $E(\theta,y)=P+y\,(\cos\theta,\sin\theta)/\max(\abs{\cos\theta},\abs{\sin\theta})$. It is
continuous and $2\pi$-periodic in $\theta$, $d(E(\theta,y),P)=y$, and $E(\theta,y)=E(\theta',y')$ if and only if $y=y'$ and $\theta'-\theta\in2\pi\Z$.
\end{enumerate}
\end{definition}

\formal{\lean{HPC.Jones.Sector.Ann}, \lean{HPC.Jones.Sector.IsCrossA}, \lean{HPC.Jones.Sector.Strip},
\lean{HPC.Jones.Sector.IsCross}, \lean{HPC.Jones.Sector.Ecov}, \lean{HPC.Jones.Sector.dist_Ecov}.}

\begin{lemma}[Lifting]\label{lem:lift}
Let $K$, $K'$ be disjoint continua crossing $\Ann$. Then $z\mapsto z-P$ has a logarithm on $K$, and $K$ has a \emph{lift}: a crossing set
$L$ of $\Strip(r_1,r_2)$ mapped by $E$ injectively onto $K$, whose translates $L+(2\pi j,0)$, $j\in\Z$, are pairwise disjoint and have union
$E^{-1}(K)$.
\end{lemma}

\begin{proof}
The open square $d(z,P)<r_1$ and the outside of the square $d(z,P)=r_2$ are connected, miss $K$, and are joined by the connected set $K'$, which
misses $K$ and meets both squares. So $P$ lies in the unbounded component of $\R^2\setminus K$. Take $y$ in that component so far away that $K$ lies in an open
half-plane missing $y$. By \vfour{Theorem~2.4}, $(z-P)/(z-y)$ has a logarithm on $K$, and $z-y$ has one (a branch of the logarithm on a half-plane); their
sum $g$ is a logarithm of $z-P$. Put $L=\set{(\Im g(z),\,d(z,P)) : z\in K}$, a continuous image of $K$, and translate it by a multiple of $2\pi$ if
desired. The properties follow from those of $E$ and from the uniqueness of logarithms up to constants in $2\pi i\Z$ on connected sets
(\vfour{Lemma~1.3}).
\end{proof}

\formal{\lean{HPC.Jones.Sector.exists_log_cross}, \lean{HPC.Jones.Sector.exists_lift} (the lift as
\lean{HPC.Jones.Sector.IsLift}), \lean{HPC.Jones.Sector.exists_normalized_lift} (a lift whose leftmost points
have first coordinate in $[0,2\pi)$).}

\section{The two sides of a crossing set}

Let $L$ be a crossing set of $\Strip=\Strip(a,b)$, $a<b$. A point of $\Strip$ is \emph{far right} (\emph{far left}) of $L$ if its first coordinate
exceeds (is less than) that of every point of $L$. The far right points form a connected set missing $L$; so do the far left points.

\begin{definition}\label{def:sides}
The \emph{right side} $\Rc(L)$ is the component of $\Strip\setminus L$ containing the far right points, the \emph{left side} $\Lc(L)$ the component
containing the far left points, and $\NR(L)=\Strip\setminus(L\cup\Rc(L))$ (the left side together with any pockets).
\end{definition}

\formal{\lean{HPC.Jones.Sector.Rc}, \lean{HPC.Jones.Sector.Lc}, \lean{HPC.Jones.Sector.NotR}, with
\lean{HPC.Jones.Sector.SideOf}, \lean{HPC.Jones.Sector.FarR}, \lean{HPC.Jones.Sector.FarL}.}

Small balls about a point of $\Strip\setminus L$, cut down to the strip, are connected and miss $L$; hence $\Rc(L)$, $\Lc(L)$ and $\NR(L)$ are open in
$\Strip$, and a connected subset of $\Strip\setminus L$ meeting one of them lies in it. Each side \emph{reaches} $L$: its closure meets $L$ (otherwise it would
be open and closed in the connected strip), so $\Rc(L)\cup L$ and $\Lc(L)\cup L$ are connected.

\formal{\lean{HPC.Jones.Sector.sideOf_absorb}, \lean{HPC.Jones.Sector.sideOf_open_rel},
\lean{HPC.Jones.Sector.GoodSide.exists_mem_closure}, \lean{HPC.Jones.Sector.GoodSide.union_isPreconnected}.}

\begin{theorem}[Strip separation]\label{thm:strip-sep}
Let $L$ be a crossing set of $\Strip(a,b)$, $a<b$. No compact connected subset of $\Strip\setminus L$ contains a point far left and a point far right of
$L$. Consequently $\Lc(L)\cap\Rc(L)=\emptyset$.
\end{theorem}

\begin{proof}
Let $\Gamma\subseteq\Strip\setminus L$ be compact and connected, containing $(\alpha,a)$ far left and $(\beta,a)$ far right of $L$. Let $p_0=(\theta_p,a)$ and $q_0=(\theta_q,b)$
be points of $L$ on the two edges; then $\alpha<\theta_p,\theta_q<\beta$. Close $\Gamma$ into
\[
Y=\Gamma\cup\bigl(\set\alpha\times[a-1,a]\bigr)\cup\bigl([\alpha,\beta]\times\set{a-1}\bigr)\cup\bigl(\set\beta\times[a-1,a]\bigr),
\]
a compact connected set. Let $p=(\theta_p,a-\tfrac12)$ and $q=(\theta_q,b+1)$. The connected set consisting of the segment from $p$ to $p_0$, the set $L$, and the
segment from $q_0$ to $q$ misses $Y$, so $p$ and $q$ lie in the same component of $\R^2\setminus Y$. By \vfour{Theorem~2.4}, $(z-p)/(z-q)$ has a logarithm on
$Y$; and $z-q$ has one, as $Y$ lies below $q$. So $z-p$ has a logarithm on $Y$; let $h$ be its imaginary part, a continuous argument of $z-p$ on $Y$.

The downward ray from $p$ meets $Y$ only in $m=(\theta_p,a-1)$, as $\Gamma$ lies in $y\ge a$. The set $Y\setminus\set m$ is connected: it is $\Gamma$ with the two
vertical segments and the two parts of the bottom segment attached. On it $h$ never takes a value in $-\frac\pi2+2\pi\Z$, so it takes values in a single interval
$(-\frac\pi2+2\pi k,\frac{3\pi}2+2\pi k)$. Approaching $m$ along the bottom segment from the left, $z-p$ points down and slightly left, so $h(z)\to\frac{3\pi}2+2\pi k$;
from the right, $h(z)\to-\frac\pi2+2\pi k$. This contradicts the continuity of $h$ at $m$.

For the consequence: if $w\in\Lc(L)\cap\Rc(L)$, the component of $w$ in $\Strip\setminus L$ contains far left and far right points of the edge $y=a$, and two points of
a component of the relatively open set $\Strip\setminus L$ are joined by a path in it, whose image is compact and connected.
\end{proof}

\formal{\lean{HPC.Jones.Sector.no_cross_join}, \lean{HPC.Jones.Sector.exists_compact_join},
\lean{HPC.Jones.Sector.lc_disjoint_rc}.}

\section{The order of disjoint crossing sets}

\begin{definition}\label{def:prec}
For disjoint crossing sets $L$, $L'$ of $\Strip(a,b)$ write $L\prec L'$ if $L'\subseteq\Rc(L)$.
\end{definition}

\formal{\lean{HPC.Jones.Sector.Prec}.}

\begin{proposition}\label{prop:prec}
On pairwise disjoint crossing sets, $\prec$ is a strict linear order. Moreover:
\begin{enumerate}[label=(\alph*)]
\item if $L\prec L'$, then $\Rc(L')\subseteq\Rc(L)$ and $\NR(L)\subseteq\NR(L')$;
\item if $L\prec L'$, then every point of $L'$ has a point of $L$ weakly to its left (first coordinate $\le$);
\item if $c>0$ and $L$ is disjoint from its translate $L+c=L+(c,0)$, then $L\prec L+c$; and $L\prec L'$ if and only if $L+c\prec L'+c$.
\end{enumerate}
\end{proposition}

\begin{proof}
A connected set missing $L$, such as $L'$, lies in $\Rc(L)$ or in $\NR(L)$.

\emph{Totality.} If $L'\subseteq\NR(L)$, then $\Rc(L)\cup L$ is a connected subset of $\Strip\setminus L'$ containing points far right of both sets, so it lies in $\Rc(L')$; thus
$L'\prec L$.

\emph{Asymmetry.} Let $L'\subseteq\Rc(L)$ and $L\subseteq\Rc(L')$. By Theorem~\ref{thm:strip-sep}, $L'$ misses $\Lc(L)$, so $\Lc(L)\cup L$ is a connected subset of
$\Strip\setminus L'$ containing far left points; it lies in $\Lc(L')$. Then $L\subseteq\Lc(L')\cap\Rc(L')=\emptyset$, which is absurd.

(a) If $L\prec L'$, then $L\not\subseteq\Rc(L')$ by asymmetry, so $L$ misses $\Rc(L')$; then $\Rc(L')\cup L'$ is a connected subset of $\Strip\setminus L$ with far right points, so
lies in $\Rc(L)$. The second inclusion follows. \emph{Transitivity:} if $L\prec L'\prec L''$, then $L''\subseteq\Rc(L')\subseteq\Rc(L)$.

(b) A point of $L'$ strictly left of all of $L$ is far left of $L$, hence in $\Lc(L)$, which is disjoint from $\Rc(L)\supseteq L'$.

(c) Translation by $c$ maps sides to sides, which gives the second statement. If $L+c\prec L$, apply (b) to a leftmost point $z$ of $L$: some point of $L+c$, whose first
coordinate is at least $z_1+c$, would be weakly left of $z$. So $L\prec L+c$ by totality.
\end{proof}

\formal{\lean{HPC.Jones.Sector.prec_total}, \lean{HPC.Jones.Sector.prec_asymm},
\lean{HPC.Jones.Sector.prec_trans}, \lean{HPC.Jones.Sector.prec_of_notR}, \lean{HPC.Jones.Sector.rc_mono},
\lean{HPC.Jones.Sector.notR_mono}, \lean{HPC.Jones.Sector.prec_min}, \lean{HPC.Jones.Sector.prec_sh_self}.}

\section{Sectors}

Let $F$ be a family of pairwise disjoint crossing sets of $\Strip(a,b)$, closed under the translations by the multiples of some $T>0$, none equal to a proper translate of
itself. By Proposition~\ref{prop:prec}(c) the translates $A+jT$ of a member $A$ increase with $j$. For $A\prec B$ in $F$ the \emph{sector} between them is
\[
\Sec(A,B)=\Rc(A)\cap\NR(B),
\]
an open subset of the strip.

\formal{\lean{HPC.Jones.Sector.LiftFam}, \lean{HPC.Jones.Sector.Sec}, \lean{HPC.Jones.Sector.sec_open_rel},
\lean{HPC.Jones.Sector.LiftFam.prec_shift_pos}.}

\begin{proposition}\label{prop:sectors}
Let $A,B,A',B',X\in F$.
\begin{enumerate}[label=(\alph*)]
\item If $A\prec X\prec B$, then $X\subseteq\Sec(A,B)$. If $X\preceq A$ or $B\preceq X$, then $X\cap\Sec(A,B)=\emptyset$.
\item If $A\prec B$, then $\Sec(A,B)\ne\emptyset$.
\item If $A\prec B\preceq A'\prec B'\preceq A+T$, then $\Sec(A,B)$ misses every translate $\Sec(A',B')+jT$.
\item If $A\prec B\prec A+T$, every point of the strip off the translates of $A$ and $B$ lies in a translate of $\Sec(A,B)$ or of $\Sec(B,A+T)$.
\end{enumerate}
\end{proposition}

\begin{proof}
(a) $X\subseteq\Rc(A)$, and $X\subseteq\NR(B)$ since $X\prec B$ (totality and asymmetry). If $X\preceq A$, then $X$ misses $\Rc(A)$; if $B\preceq X$, then $X$ is $B$ or lies in
$\Rc(B)$, and misses $\NR(B)$.

(b) A point $x\in A$ lies in the closure of $\Rc(A)$ (the side reaches $A$), and $x\in\NR(B)$ by (a). As $\NR(B)$ is open in the strip, points of $\Rc(A)$ near $x$ lie in
$\Sec(A,B)$.

(c) For $j\ge0$: $\Sec(A',B')+jT\subseteq\Rc(A'+jT)\subseteq\Rc(B)$, since $B\preceq A'\preceq A'+jT$, while $\Sec(A,B)\subseteq\NR(B)$. For $j<0$: $\Sec(A',B')+jT\subseteq\NR(B'+jT)\subseteq\NR(A)$,
since $B'+jT\preceq A+(j+1)T\preceq A$, while $\Sec(A,B)\subseteq\Rc(A)$.

(d) Let $y$ be such a point. For $j$ very negative $y$ is far right of $A+jT$, and for $j$ very large it is far left; and $\Rc(A+jT)$ decreases with $j$. Let $j$ be largest with $y\in
\Rc(A+jT)$; then $y\in\NR(A+(j+1)T)$. If $y\in\Rc(B+jT)$, then $y\in\Sec(B+jT,A+(j+1)T)=\Sec(B,A+T)+jT$; otherwise $y\in\NR(B+jT)$ and $y\in\Sec(A,B)+jT$.
\end{proof}

\formal{\lean{HPC.Jones.Sector.LiftFam.sub_sec}, \lean{HPC.Jones.Sector.LiftFam.sec_disjoint_mem},
\lean{HPC.Jones.Sector.LiftFam.sec_nonempty}, \lean{HPC.Jones.Sector.LiftFam.sec_disjoint},
\lean{HPC.Jones.Sector.LiftFam.sec_cover}.}

\section{The sector theorem and the cyclic order}

\begin{theorem}[The sector theorem]\label{thm:sector}
Let $K$, $K'$ be disjoint continua crossing $\Ann$. Then $\Ann\setminus(K\cup K')=V_1\cup V_2$ with $V_1$, $V_2$ nonempty, disjoint and open in $\Ann$. Consequently a connected
subset of $\Ann$ meeting $V_1$ and $V_2$ meets $K\cup K'$.
\end{theorem}

\begin{proof}
Take lifts $A$ of $K$ and $B$ of $K'$ (Lemma~\ref{lem:lift}); all translates of $A$ and $B$ form a family $F$ as above, with $T=2\pi$. Translating $B$, we may assume $A\prec B\prec
A+2\pi$. Let $V_1=E(\Sec(A,B))$ and $V_2=E(\Sec(B,A+2\pi))$. They are nonempty by Proposition~\ref{prop:sectors}(b); disjoint by (c), because $E(u)=E(v)$ only if $v$ is a translate of $u$
by a multiple of $2\pi$; and they cover $\Ann\setminus(K\cup K')$ by (d), since $E^{-1}(K\cup K')$ is the union of the translates of $A$ and $B$. They are open in $\Ann$ because the sectors
are open in the strip and $E$ maps small relatively open sets of the strip homeomorphically onto relatively open sets of $\Ann$. The consequence is immediate.
\end{proof}

\formal{\lean{HPC.Jones.sector_theorem}, from \lean{HPC.Jones.Sector.sector_two} and
\lean{HPC.Jones.Sector.two_piece_lift}; Moore uses it on p.~236 of his 1922 paper (``must clearly contain a
point either of $k_{n-1}$ or of $k_{n+1}$'').}

\begin{theorem}[The cyclic order]\label{thm:cyclic}
Let $k_0,k_1,\dots$ be pairwise disjoint continua crossing $\Ann$. There are a subsequence $k_{g(0)},k_{g(1)},\dots$ and, for $n\ge1$, sets $V_n\subseteq\Ann$ such that
\begin{enumerate}[label=(\roman*)]
\item $k_{g(n)}\subseteq V_n$, and $V_n$ misses $k_{g(m)}$ for $m\ne n$;
\item $V_n\cap V_m=\emptyset$ if $\abs{n-m}\ge2$;
\item a connected subset of $\Ann\setminus(k_{g(n-1)}\cup k_{g(n+1)})$ that meets $V_n$ lies in $V_n$.
\end{enumerate}
\end{theorem}

\begin{proof}
Choose normalized lifts $L_n$ of the $k_n$: the least first coordinate of $L_n$ lies in $[0,2\pi)$. By Proposition~\ref{prop:prec}(b), $L_m\prec L_n+2\pi$ for all $m,n$. Any sequence in a linear
order has a strictly monotone subsequence; reversing the orientation if necessary, let $L_{g(0)}\prec L_{g(1)}\prec\cdots$. For $n\ge1$ put $V_n=E\bigl(\Sec(L_{g(n-1)},L_{g(n+1)})\bigr)$.
Then (i) follows from Proposition~\ref{prop:sectors}(a) (applied to all translates, using $L_m\prec L_n+2\pi$), (ii) from Proposition~\ref{prop:sectors}(c), and (iii) from the proof of
Theorem~\ref{thm:sector}: $V_n$ is one of the two pieces cut off by $k_{g(n-1)}$ and $k_{g(n+1)}$, and a connected set in the union of two disjoint relatively open pieces lies in one of
them.
\end{proof}

\formal{\lean{HPC.Jones.sector_cyclic_order}, from \lean{HPC.Jones.Sector.sector_cyclic},
\lean{HPC.Jones.Sector.sector_props}, \lean{HPC.Jones.Sector.preconnected_sub_piece}. Moore (p.~237) orders the
continua by their points on a middle circle; the order of the lifts avoids having to show that this is well
defined.}

\section{Crossing continua from a failure of local connectedness}

\begin{definition}\label{def:cik}
A set $M\subseteq\R^2$ is \emph{connected im kleinen at $P\in M$} if for every $\varepsilon>0$ there is $\delta>0$ such that every point of $M$ within $\delta$ of $P$ lies with $P$ in a connected
subset of $M$ of diameter less than $\varepsilon$.
\end{definition}

\formal{\lean{HPC.Jones.CikAt}.}

\begin{proposition}\label{prop:cik}
\leavevmode
\begin{enumerate}[label=(\alph*)]
\item A set $M\subseteq\R^2$ connected im kleinen at each of its points is locally connected.
\item If a closed connected set $M\subseteq\R^2$ is not connected im kleinen at $P\in M$, there are $0<r_1<r_2$ and infinitely many pairwise disjoint continua in $M$ crossing $\Ann(P,r_1,r_2)$.
\end{enumerate}
\end{proposition}

\begin{proof}
(a) Let $F\subseteq M$ be open and $y$ a point of the component $C$ of $y_0$ in $F$; take $\varepsilon$ with $B(y,\varepsilon)\cap M\subseteq F$ and $\delta$ for $\varepsilon$ at $y$. A point $z$ of $M$ within $\delta$ of $y$
lies with $y$ in a connected subset of $M$ of diameter less than $\varepsilon$, which lies in $F$, hence in $C$. So components of open sets are open.

(b) Let $\varepsilon$ witness the failure, $r_2=\varepsilon/3$, $r_1=\varepsilon/6$, and let $\Gamma$ be the component of $P$ in the compact set $M\cap\bar B(P,r_2)$. Points of $M$ arbitrarily close to $P$ lie
outside $\Gamma$ (otherwise $\Gamma$, of diameter less than $\varepsilon$, would do). Their components in $M\cap\bar B(P,r_2)$ are infinitely many: a finite union of components missing $P$ is closed
and misses a neighbourhood of $P$. They are pairwise disjoint, and each meets the outer square by boundary bumping (\vthree{Theorem~1.10}), as $M\not\subseteq\bar B(P,r_2)$ (otherwise $M$ itself, of diameter less than $\varepsilon$, would do). Infinitely many of them
contain points within $r_1$ of $P$. In each of these, a second application of boundary bumping, from a point on the outer square in the part outside $B(P,r_1)$, gives a continuum crossing the
annulus.
\end{proof}

\formal{\lean{HPC.Jones.locallyConnectedSpace_of_cikAt}, \lean{HPC.Jones.exists_crossing_continua}, and for closed
sets \lean{HPC.Jones.Sector.exists_crossing_continua_closed}. Moore attributes (b) to ``the argument beginning
with the second line from the bottom of page 365'' of a 1918 note.}

\section*{Exercises}

\begin{exercise}
Show that Theorem~\ref{thm:sector} fails if $K$ and $K'$ are allowed to meet, and also if $K$ is only required to meet one of the two squares.
\end{exercise}

\begin{exercise}
Prove Theorem~\ref{thm:strip-sep} for a crossing set $L$ which is a vertical segment directly, and explain why the general case cannot be reduced to it.
\end{exercise}

\begin{exercise}
Let $k_n$ be the radial segments of $\Ann(0,1,2)$ at the angles $1/n$, $n\ge1$. Describe the subsequence and the sets $V_n$ of Theorem~\ref{thm:cyclic}.
\end{exercise}

\begin{exercise}
Show that the $\sin(1/x)$-continuum is not connected im kleinen at the points of its limit segment, and find the crossing continua of Proposition~\ref{prop:cik}(b).
\end{exercise}

%======================================================================
\chapter{Boundaries of complementary domains}\label{ch:boundaries}
%======================================================================

Jones's 1949 proof continues: ``Hence each element of $G$ is a simple domain [Whyburn]. Let the simple closed
curve $J$ denote the boundary of an element $D$ of $G$.'' Here $G$ is the collection of complementary domains
of $M$, and the reference is to Whyburn's 1939 paper \emph{Semi-locally-connected sets}, \S14:

\begin{quote}
\emph{Theorem~14.} If $M$ is a compact semi-locally-connected plane continuum, the boundary of each complementary
domain of $M$ is locally connected.

\emph{Corollary~1.} If $M$ is moreover cyclic (has no cut point), the boundary of every complementary domain of $M$
is a simple closed curve.
\end{quote}

Whyburn's proof of Theorem~14 reduces to bounded domains by an inversion, joins points by arcs and
uses a simple closed curve supplied by ``a theorem of R.~L.~Moore'' in one step, and a curve built
from a point of order~2 in another. The reduction has a gap: the centre of the inversion must lie in a
second complementary domain (a centre in $R$ leaves the image of $R$ unbounded, and a centre on $M$ destroys
compactness), and when $R$ is the only complementary domain, for instance when $M$ is a closed disk, there is none. The proof below needs no inversion, no arcs and no Jordan curves.
It uses Janiszewski's theorem (\vfour{Theorem~2.7}), the sector theorem of Chapter~\ref{ch:sectors}, and
Whyburn's theorem on cut points of order other than~$2$ (Corollary~\ref{cor:whyburn}(c)). Corollary~1, which
Whyburn states without proof, then follows from the arc theorem for Peano continua and a lemma on
$\theta$-curves.

Throughout, $M\subseteq\R^2$ is a compact continuum that is semi-locally connected at each of its points
(Definition~\ref{def:slc}), $R$ is a component of $\R^2\setminus M$, and $B=\fr R$. $B$ is a continuum by Brouwer's
theorem (\vfour{Theorem~1.12}), and $B\subseteq M$. Distances are taken in the maximum metric, as in
Chapter~\ref{ch:sectors}; $V_\varepsilon(p)$ is the open ball $B(p,\varepsilon)$.

\formal{\lean{HPC.Jones.isConnected_frontier_complDomain}.}

\section{Joining points of a domain near a non-cut point}

\begin{lemma}[Whyburn (6.22), the form used in \S14]\label{lem:whyburn622}
Let $X$ be a continuum semi-locally connected at each point, and $p$ a non-cut point of $X$. For every $\varepsilon>0$ there is $\delta\in(0,\varepsilon)$
such that some component of $X\setminus V_\delta(p)$ contains $X\setminus V_\varepsilon(p)$.
\end{lemma}

\begin{proof}
If $X=\set p$ there is nothing to prove; fix $q\ne p$. By Corollary~\ref{cor:jones41}, $X$ is aposyndetic. For each $a\ne p$ there is a continuum $H_a$
missing $p$ and containing an open $U_a\ni a$; by Proposition~\ref{prop:whyburn62} (the component of $a$ in the connected set $X\setminus\set p$ is all of it), some
continuum missing $p$ contains $a$ and $q$; so we may assume $q\in H_a$. Finitely many $U_a$ cover the compact set $X\setminus V_\varepsilon(p)$. The union $W$ of the
corresponding $H_a$ is a continuum (all contain $q$) missing $p$; choose $\delta<\varepsilon$ with $V_\delta(p)\cap W=\emptyset$. Then $W$ is a connected subset of $X\setminus V_\delta(p)$
containing $X\setminus V_\varepsilon(p)$.
\end{proof}

\formal{\lean{HPC.Jones.whyburn39_6_22_sec14}.}

\begin{lemma}[The joining lemma]\label{lem:join}
Let $N\subseteq M$ be a continuum containing $M\setminus V_\varepsilon(p)$, and let $W$ be the component of $\R^2\setminus\bigl(N\cup(\R^2\setminus V_\varepsilon(p))\bigr)$ containing $p$. Any two
points $a,b\in R\cap W$ lie in the same component of $W\setminus M$.
\end{lemma}

\begin{proof}
Let $K=(M\cap\cl W)\cup N$ and $L=(\R^2\setminus W)\cap\bar B(p,\rho)$, with $\rho>\varepsilon$ so large that $M\subseteq\bar B(p,\rho)$. Both are compact, and $K\cap L=N$ (a point of $M\cap\cl W$ outside $W$ lies in the frontier of $W$,
hence in $N$ or outside $V_\varepsilon(p)$, and in the latter case in $M\setminus V_\varepsilon(p)\subseteq N$), which is connected. $K\subseteq M$ does not separate $a$ from $b$, as both lie in the component
$R$ of $\R^2\setminus M$; and $L$ does not separate them, as both lie in the connected set $W$, which misses $L$. By Janiszewski's theorem (\vfour{Theorem~2.7}), $K\cup L$ does not
separate $a$ from $b$. The component $C$ of $a$ in $\R^2\setminus(K\cup L)$ is a connected set starting in $W$; it cannot leave $W$, because the frontier of $W$ lies in $N\subseteq K$ or in
$\cl{V_\varepsilon(p)}\setminus V_\varepsilon(p)\subseteq L$. So $C\subseteq W$, and $C$ misses $M\cap\cl W\subseteq K$: $C\subseteq W\setminus M$, and $b\in C$.
\end{proof}

\formal{\lean{HPC.Jones.Sec14.join_in_domain}; it replaces Whyburn's $\theta$-curve step and Moore's curve
$C$ (\lean{HPC.Jones.whyburn39_moore_curve} records the form used).}

\begin{proposition}[Step 3 of Whyburn's proof]\label{prop:local-join}
Let $p\in B$ be a non-cut point of $M$. For every $\varepsilon>0$ there is $\delta>0$ such that any two points of $R$ within $\delta$ of $p$ lie in a connected subset of $R$ within
$\varepsilon$ of $p$. In particular there is no sequence of points of $R$ converging to $p$ no two of which lie in a connected subset of $R$ of diameter less than $\varepsilon$.
\end{proposition}

\begin{proof}
Lemma~\ref{lem:whyburn622} gives $\delta<\varepsilon$ and a component of $M\setminus V_\delta(p)$ containing $M\setminus V_\varepsilon(p)$; let $N$ be its closure, a continuum in $M\setminus V_\delta(p)$. The ball
$V_\delta(p)$ is connected, misses $N$ and lies in $V_\varepsilon(p)$, so it lies in the component $W$ of Lemma~\ref{lem:join}. For $a,b\in R\cap V_\delta(p)$, the component of $a$ in $W\setminus M$ is a
connected subset of $R$ (it is connected, misses $M$ and contains $a$) inside $W\subseteq V_\varepsilon(p)$, containing $b$. For the second statement apply the first with $\varepsilon/3$.
\end{proof}

\formal{\lean{HPC.Jones.Sec14.local_join_of_noncut}, \lean{HPC.Jones.whyburn39_local_noncut}.}

\section{Theorem 14}

Say that Whyburn's condition \emph{fails at $p\in B$ with constant $\varepsilon$} if some sequence of points of $R$ converges to $p$ and no two of its terms lie in a connected subset of
$R$ of diameter less than $\varepsilon$.

\formal{\lean{HPC.Jones.Sec14.BadAt}.}

\begin{lemma}[Accumulating continua exclude finite order]\label{lem:accum}
Let $T_0,T_1,\dots$ be pairwise disjoint connected subsets of a metric space, $b_i\in T_i$ with $b_i\to p$, and suppose each $T_i$ has a point at distance at least $d>0$ from $p$. Then $p$
has no finite order.
\end{lemma}

\begin{proof}
Let $V\ni p$ be open inside $B(p,d)$. For all large $i$, $b_i\in V$, and $T_i$ has a point outside $V$; so the connected set $T_i$ meets $\fr V$. As the $T_i$ are disjoint, $\fr V$ is infinite.
\end{proof}

\formal{\lean{HPC.Jones.Sec14.not_orderLE_of_accumulating}, recorded as
\lean{HPC.Jones.whyburn39_curve_from_order_two}; it replaces Whyburn's curve $K$ and his pigeonhole argument.}

\begin{lemma}[Bad points at every middle radius]\label{lem:bad-point}
Let $t_0,t_1,\dots$ be pairwise disjoint continua in $B$ crossing $\Ann(P,r_1,r_2)$, and let $\rho\in[r_1+\frac{r_2-r_1}3,\,r_2-\frac{r_2-r_1}3]$. There is $h\in B$ with $d(h,P)=\rho$ at which Whyburn's
condition fails with constant $\eta=\frac{r_2-r_1}6$, and at which infinitely many of the $t_n$ accumulate (points $b_i\in t_{s(i)}$, $s$ strictly increasing, with $b_i\to h$).
\end{lemma}

\begin{proof}
Apply the cyclic order (Theorem~\ref{thm:cyclic}) to the $t_n$; renumber the subsequence as $t_0,t_1,\dots$, with the sets $V_n$. Each $t_n$ crosses the annulus, so it has a point $b_n$ with
$d(b_n,P)=\rho$. The points $b_n$ with $n$ odd have a limit point $h$ (the square $d(z,P)=\rho$ is compact), $h\in B$ as $B$ is closed, and $d(h,P)=\rho$. For each odd $n\ge1$ choose $\delta_n\le\min(\eta,1/n)$
such that the ball $B(b_n,\delta_n)$ misses $t_{n-1}\cup t_{n+1}$ and lies, within $\Ann$, in $V_n$; since $b_n\in B=\fr R$, there is $c_n\in R\cap B(b_n,\delta_n)$. Along the odd indices for which
$b_n\to h$, also $c_n\to h$.

Let $n<m$ be two of these indices and $C\subseteq R$ a connected set of diameter less than $\eta$ containing $c_n$ and $c_m$. $C$ misses every $t_k\subseteq B$, since $R$ is open and $B=\fr R$. As $d(c_n,P)$ is
within $\eta$ of $\rho$, every point of $C$ is within $2\eta=\frac{r_2-r_1}3$ of the square $d(z,P)=\rho$, so $C$ lies in $\Ann$. It meets $V_n$ (at $c_n$) and misses $t_{n-1}\cup t_{n+1}$, so $C\subseteq V_n$ by Theorem~\ref{thm:cyclic}(iii). But $c_m\in V_m$ and $\abs{n-m}\ge2$, so $V_n\cap V_m=\emptyset$
by Theorem~\ref{thm:cyclic}(ii). So Whyburn's condition fails at $h$.
\end{proof}

\formal{\lean{HPC.Jones.Sec14.exists_bad_point_of_crossing}. This is Moore's 1922 argument (p.~236) at an
arbitrary radius; at the middle radius it gives Moore's local sufficiency theorem
\lean{HPC.Jones.moore22_sufficiency_local}, \lean{HPC.Jones.Sector.sufficiency_local}.}

\begin{theorem}[Whyburn 1939, Theorem~14]\label{thm:whyburn14}
Let $M\subseteq\R^2$ be a compact continuum, semi-locally connected at each point. The boundary $B$ of each complementary domain $R$ of $M$, bounded or not, is locally connected.
\end{theorem}

\begin{proof}
Suppose not. By Proposition~\ref{prop:cik}(a), $B$ is not connected im kleinen at some $P\in B$, and by Proposition~\ref{prop:cik}(b) there are pairwise disjoint continua $t_n\subseteq B$ crossing an
annulus $\Ann(P,r_1,r_2)$. For each $\rho$ in the middle third $I_0=[r_1+\frac{r_2-r_1}3,r_2-\frac{r_2-r_1}3]$, Lemma~\ref{lem:bad-point} gives $h_\rho\in B\subseteq M$ with $d(h_\rho,P)=\rho$. At $h_\rho$ infinitely many
disjoint continua $t_{s(i)}\subseteq M$ accumulate, each reaching distance at least $\frac{r_2-r_1}3$ from $h_\rho$; by Lemma~\ref{lem:accum}, $h_\rho$ has no finite order in $M$, in particular not order~$2$.
Whyburn's condition fails at $h_\rho$, so by Proposition~\ref{prop:local-join}, $h_\rho$ is a cut point of $M$. The points $h_\rho$, $\rho\in I_0$, are distinct. So $M$ has uncountably many cut points not of
order~$2$, against Corollary~\ref{cor:whyburn}(c).
\end{proof}

\formal{\lean{HPC.Jones.whyburn39_theorem14}.}

\section{Corollary 1}

Recall that an \emph{arc} is a homeomorphic image of $[0,1]$ (\vfive{Definition~1.1}). A \emph{$\theta$-curve} consists of three arcs $K_1,K_2,K_3$ from $a$ to $b$, $a\ne b$, any two of which meet only in
$\set{a,b}$. Then each $J_{ij}=K_i\cup K_j$ is a simple closed curve, and $\oarc K_k=K_k\setminus\set{a,b}$ is a connected set missing $J_{ij}$, for $\set{i,j,k}=\set{1,2,3}$.

\begin{lemma}[The $\theta$-curve lemma]\label{lem:theta}
Let $K_1\cup K_2\cup K_3\subseteq\R^2$ be a $\theta$-curve and $k_i\in\oarc K_i$.
\begin{enumerate}[label=(\alph*)]
\item \emph{(No common side.)} No point $m$ off the $\theta$-curve lies, for each $i$, in the component of $\R^2\setminus J_{jk}$ containing $k_i$.
\item No connected set $D$ missing the $\theta$-curve has the whole $\theta$-curve in its closure.
\end{enumerate}
\end{lemma}

\begin{proof}
We use the Jordan curve theorem (\vfive{Theorem~4.8}): $\R^2\setminus J_{ij}$ has exactly two components, each with frontier $J_{ij}$. Hence a connected set $A$ missing $J_{ij}$ whose closure contains a
connected set $A'$ missing $J_{ij}$ lies in the same component as $A'$: if $A$ lies in the component $S$, then $\cl A\subseteq S\cup J_{ij}$, so $A'\subseteq S$.

(a) Let $S$ be the component of $\R^2\setminus J_{12}$ containing $\oarc K_3$ (and $m$), and $T$ the other one; pick $w\in T$. $T$ is connected and misses $K_3$ (as $\oarc K_3\subseteq S$), so it misses $J_{13}$ and $J_{23}$. Its
closure contains $\fr T=J_{12}\supseteq\oarc K_2$, so $T$ lies in the component of $\R^2\setminus J_{13}$ containing $\oarc K_2$, which also contains $m$. Likewise $T$ and $m$ lie in the same component of
$\R^2\setminus J_{23}$. So neither $J_{13}$ nor $J_{23}$ separates $m$ from $w$, and $J_{13}\cap J_{23}=K_3$ is connected. By Janiszewski's theorem, $J_{13}\cup J_{23}$ does not separate $m$ from $w$. But it contains
$J_{12}$, which does.

(b) Let $m\in D$. For each $i$, $D$ is a connected set missing $J_{jk}$ whose closure contains $\oarc K_i$; so $D$, and $m$, lie in the component of $\R^2\setminus J_{jk}$ containing $k_i$. This contradicts (a).
\end{proof}

\formal{\lean{HPC.Jones.theta_no_common_side}, \lean{HPC.Jones.theta_not_subset_closure} (with
\lean{HPC.Jones.jordan_side}).}

\begin{lemma}[The boundary is cyclic]\label{lem:frontier-cyclic}
If $z$ is a cut point of $B$, it is a cut point of $M$.
\end{lemma}

\begin{proof}
Let $B\setminus\set z=U\cup V$ with $U$, $V$ separated and nonempty, and put $P=U\cup\set z$, $Q=V\cup\set z$: compact sets with $P\cap Q=\set z$. Let $R_P$, $R_Q$ be the components of $\R^2\setminus P$ and
$\R^2\setminus Q$ containing the connected set $R$. A point of $R_P\cap R_Q$ is separated from $R$ neither by $P$ nor by $Q$, so by Janiszewski's theorem not by $P\cup Q=B$; as $R$ is a component of $\R^2\setminus B$,
$R_P\cap R_Q=R$. Each $v\in V$ is a limit of points of $R\subseteq R_P$ and lies outside $P$, so $V\subseteq R_P$ ($R_P$ is closed in $\R^2\setminus P$); likewise $U\subseteq R_Q$.

Consider $M\cap R_Q$, an open subset of $M\setminus\set z$ containing $U$ and missing $V\subseteq Q$. It is also closed in $M\setminus\set z$: a point of its closure lies in $R_Q\cup Q$; if it is a point $v\in V$, a small
connected ball about $v$ missing $P$ lies in $R_P$, so its intersection with $R_Q$ lies in $R_P\cap R_Q=R$, which misses $M$, a contradiction. So $M\setminus\set z$ is not connected.
\end{proof}

\formal{\lean{HPC.Jones.frontier_cyclic}, \lean{HPC.Jones.frontier_cyclic'}.}

\begin{proposition}[Curves in cyclic Peano continua]\label{prop:cyclic-peano}
Let $X$ be a locally connected continuum.
\begin{enumerate}[label=(\alph*)]
\item Two distinct points of a connected open subset $U$ of $X$ are the end points of an arc in $U$.
\item If $X$ has no cut point and is nondegenerate, it contains a simple closed curve; if moreover $X$ is not a simple closed curve, it contains a $\theta$-curve.
\end{enumerate}
\end{proposition}

\begin{proof}
(a) Moore's nested simple chains, with the connected open subsets of $U$ as regions, give a compact connected subset of $U$ irreducible between the two points (as in \vfive{Theorem~5.7}); it is an arc
by \vfive{Corollary~1.9}.

(b) Take an arc $\alpha$ from $a$ to $b$ and a point $c$ of $\alpha$ other than $a,b$. As $X\setminus\set c$ is a connected open set, (a) gives an arc $\beta$ from $a$ to $b$ avoiding $c$. Let $a_1$ be the last point of $\beta$
on the sub-arc of $\alpha$ from $a$ to $c$, and $b_1$ the first point of $\beta$ after $a_1$ on $\alpha$; then $b_1$ lies beyond $c$, and the sub-arcs of $\alpha$ and of $\beta$ from $a_1$ to $b_1$ form a simple closed curve
$J$. If $X\ne J$, take $w\notin J$, an arc from $w$ to $J$ and its first point $c_1$ on $J$, giving an arc $\gamma_1$ from $w$ to $c_1$ meeting $J$ only in $c_1$. In $X\setminus\set{c_1}$ take an arc $\delta$ from $w$ to a
point of $J$; let $c_2$ be its first point on $J$ and $u$ its last point on $\gamma_1$ before $c_2$. Then the sub-arc of $\gamma_1$ from $c_1$ to $u$ followed by the sub-arc of $\delta$ from $u$ to $c_2$ is an arc from $c_1$ to $c_2\ne
c_1$ meeting $J$ only in its end points; with the two arcs of $J$ between $c_1$ and $c_2$ it forms a $\theta$-curve.
\end{proof}

\formal{\lean{HPC.Jones.exists_tarc_in_open}, \lean{HPC.Jones.exists_tarc_of_locallyConnected},
\lean{HPC.Jones.exists_scc_of_cyclic}, \lean{HPC.Jones.exists_theta_of_cyclic},
\lean{HPC.Jones.theta_of_cyclic_peano}.}

\begin{theorem}[Whyburn 1939, \S14, Corollary~1]\label{thm:whyburn-cor1}
Let $M\subseteq\R^2$ be a nondegenerate compact continuum, semi-locally connected at each point and without cut points. The boundary $B$ of each complementary domain $R$ of $M$ is a simple closed curve.
\end{theorem}

\begin{proof}
$B$ is a continuum, locally connected by Theorem~\ref{thm:whyburn14}, and without cut points by Lemma~\ref{lem:frontier-cyclic}. It is nondegenerate: if $B\subseteq\set z$, the connected set $\R^2\setminus\set z$ would lie in the union of the disjoint open sets $R$ and
$\R^2\setminus\cl R$, hence in $R$, and $M\subseteq\set z$. If $B$ is not a simple closed curve, it contains a $\theta$-curve (Proposition~\ref{prop:cyclic-peano}(b)). This $\theta$-curve lies in $B\subseteq\cl R$, and $R$ is connected and misses it, against
Lemma~\ref{lem:theta}(b). So $B$ is a simple closed curve, homeomorphic to the circle (\vfive{Proposition~1.13}).
\end{proof}

\formal{\lean{HPC.Jones.whyburn39_corollary1}, with \lean{HPC.Jones.tscc_homeomorph_circle}. Whyburn gives no
proof and refers to Wilder (1938). The hypothesis ``no cut points'' is stated there as: $M\setminus\set z$ is connected and nonempty for
every $z\in M$, which includes nondegeneracy. The same argument also shows that every point of $B$ is accessible from $R$:
\lean{HPC.Jones.accessible_of_cyclic}, which the classification does not need.}

\section*{Exercises}

\begin{exercise}
Show that Theorem~\ref{thm:whyburn14} fails without semi-local connectedness: find a plane continuum one of whose complementary domains has a boundary that is not locally connected.
\end{exercise}

\begin{exercise}
Let $M$ be the union of two circles meeting in one point. Determine the boundaries of the complementary domains, and explain which hypothesis of Theorem~\ref{thm:whyburn-cor1} fails.
\end{exercise}

\begin{exercise}
Prove Lemma~\ref{lem:theta}(b) directly for the $\theta$-curve consisting of a circle and a diameter.
\end{exercise}

%======================================================================
\chapter{Homogeneous aposyndetic plane continua are circles}\label{ch:jones49}
%======================================================================

This chapter proves the theorem of Jones's \emph{Note on homogeneous plane continua} (1949), the last
ingredient of the decomposable case:

\begin{quote}
\emph{Theorem~1.} If a compact plane continuum $M$ is both homogeneous and aposyndetic, then $M$ is a simple
closed curve.
\end{quote}

As printed the theorem is false for a one-point continuum, which is homogeneous and aposyndetic; it is proved
for nondegenerate $M$. Jones's proof occupies about a page. It is followed here step by step, with three
additions: the local separating points at the end of both cases are produced explicitly, the claim ``$M-(A+B)$ is
not connected'' in Case~2 is proved with the $\theta$-curve lemma, and in Case~1 the Baire argument ends with a
local separating point rather than with Jones's remark on limit points.

Throughout, $M\subseteq\R^2$ is a nondegenerate compact continuum, $G$ is the collection of components of $\R^2\setminus M$
(its \emph{complementary domains}), and for $x\notin M$, $\comp Mx$ is the complementary domain containing $x$.

\formal{The statement is \lean{HPC.Jones.Jones49Thm1}; the one-point counterexample to the printed version is
\lean{HPC.Jones.not_jones49_theorem1_point}.}

\section{The setting}

\begin{lemma}[Jones 1949, p.~113]\label{lem:jones49-setting}
Let $M$ be homogeneous and aposyndetic. Then $M$ has no cut point, $M$ is semi-locally connected at each point, and the boundary of each complementary
domain of $M$ is a simple closed curve.
\end{lemma}

\begin{proof}
A nondegenerate continuum has a non-cut point (\vthree{Theorem~3.6}); a homeomorphism of $M$ carries it to any given point, and homeomorphisms carry
non-cut points to non-cut points. Semi-local connectedness is Corollary~\ref{cor:jones41}. The last statement is Theorem~\ref{thm:whyburn-cor1}.
\end{proof}

\formal{\lean{HPC.Jones.jones49_setting}, \lean{HPC.Jones.jones49_simple_domains}. Jones: ``Because $M$ is homogeneous and
contains a non-separating point, no point of $M$ separates $M$. Since $M$ is aposyndetic, $M$ is semi-locally-connected. Hence each
element of $G$ is a simple domain.''}

\begin{lemma}[Two sources of local separating points]\label{lem:local-sep}
\leavevmode
\begin{enumerate}[label=(\alph*)]
\item If a point $p\in M$ has a neighbourhood in $M$ contained in a simple closed curve $J\subseteq M$, then $p$ is a local separating point of $M$.
\item If $a\ne b$ are points of $M$, $M\setminus\set b$ is connected, and $M\setminus\set{a,b}=X_1\cup X_2$ with $X_1$, $X_2$ nonempty and separated, then $a$ is a local
separating point of $M$.
\end{enumerate}
\end{lemma}

\begin{proof}
(a) Let $W\ni p$ be open in $M$ with $W\subseteq J$, and $q\in J\setminus\set p$. The points $p$, $q$ split $J$ into two arcs $L_1$, $L_2$. The sets $W\cap\oarc L_1$ and $W\cap\oarc L_2$ are open in $J$, hence in $W$ and in $M$;
they cover $(W\setminus\set q)\setminus\set p$ and are disjoint; and small sub-arcs of $L_1$ and $L_2$ at $p$ are connected subsets of $W\setminus\set q$ through $p$ meeting both. So $p$ separates the open set
$W\setminus\set q$.

(b) $a$ separates the open set $R=M\setminus\set b$, whose component containing $a$ is all of $R$: take for $A$, $B$ the complements of $\cl{X_2}$ and $\cl{X_1}$.
\end{proof}

\formal{\lean{HPC.Jones.isLocalSepPoint_of_open_sub_tscc}, \lean{HPC.Jones.isLocalSepPoint_of_separation}.}

From now on assume that $M$ is homogeneous and aposyndetic and that \emph{no point of $M$ has order~$2$}; we derive a contradiction. Fix $x\notin M$, let
$D=\comp Mx$, and let $J=\fr D$, a simple closed curve by Lemma~\ref{lem:jones49-setting}. By Corollary~\ref{cor:some-order-two} (a nondegenerate continuum is
uncountable), \emph{no point of $M$ is a local separating point}.

\section{Case 1: \texorpdfstring{$M\setminus J$}{M - J} is connected}

\begin{lemma}\label{lem:frontier-of-jordan}
Let $J'\subseteq M$ be a simple closed curve and $A$ a connected set missing $J'$ with $M\subseteq A\cup J'$. Then $J'$ is the boundary of a complementary domain of $M$.
\end{lemma}

\begin{proof}
By the Jordan curve theorem (\vfive{Theorem~4.8}) $A$ lies in one component of $\R^2\setminus J'$; let $T$ be the other one. $T$ misses $M\subseteq A\cup J'$, is open and connected, and
its frontier is $J'\subseteq M$; so $T$ is a component of $\R^2\setminus M$, with boundary $J'$.
\end{proof}

\formal{\lean{HPC.Jones.frontier_cc_of_jordan}.}

\begin{proposition}[Case 1]\label{prop:case1}
If $M\setminus J$ is connected, some point of $M$ is a local separating point.
\end{proposition}

\begin{proof}
\emph{Every point of $M$ lies on the boundary of a complementary domain.} Given $p\in M$, let $h$ be a homeomorphism of $M$ carrying a point of $J$ to $p$. Then $J'=h(J)\ni p$ is a simple
closed curve in $M$ and $M\setminus J'=h(M\setminus J)$ is connected; apply Lemma~\ref{lem:frontier-of-jordan} with $A=M\setminus J'$.

\emph{$G$ is countable}, being a family of disjoint nonempty open sets in the plane (Lemma~\ref{lem:disjoint-open-countable}).

So $M$ is the union of countably many closed sets $\fr U$, $U\in G$. By the Baire category theorem one of them, $J_i$, contains a nonempty open subset $W$ of $M$. $J_i$ is a simple closed curve
(Lemma~\ref{lem:jones49-setting}), so every point of $W$ is a local separating point by Lemma~\ref{lem:local-sep}(a).
\end{proof}

\formal{\lean{HPC.Jones.exists_frontier_through}, \lean{HPC.Jones.countable_complDomains},
\lean{HPC.Jones.jones49_case1_core}, \lean{HPC.Jones.jones49_case1}. Jones ends with: ``Since no point of $M$ is of order~2 in
$M$, each point of $J_i$ is a limit point of $M-J_i$. This contradicts a well known theorem (Baire).''}

\section{Case 2: \texorpdfstring{$M\setminus J$}{M - J} is not connected}

\begin{lemma}[The separation of $M\setminus\set{A,B}$]\label{lem:theta-sep}
Let $J_1$, $J_2$, $L$ be arcs in $M$ from $A$ to $B$ forming a $\theta$-curve. Suppose that $D$ is a complementary domain of $M$ with boundary $J_1\cup J_2$, and $U\ne D$ a complementary domain of
$M$ whose boundary is a simple closed curve $L\cup L'$, where $L'$ is an arc from $A$ to $B$ with $L\cap L'=\set{A,B}$. Then $M\setminus\set{A,B}$ is the union of two nonempty separated sets.
\end{lemma}

\begin{proof}
Let $E_1$ be the side of $J_1\cup L$ not containing $\oarc J_2$, and $E_2$ the side of $J_2\cup L$ not containing $\oarc J_1$. Since $D$
misses $\oarc L\subseteq M$ and has frontier $J_1\cup J_2$, $D$ is the side of $J_1\cup J_2$ not containing $\oarc L$.

\emph{$D$, $E_1$, $E_2$ are pairwise disjoint and cover the complement of the $\theta$-curve.} $E_1$ misses $J_2$ (it misses $\oarc J_2$ and $A,B$), so it is a
connected set missing $J_2\cup L$ and $J_1\cup J_2$, with $\oarc J_1$ and $\oarc L$ in its closure; by the remark at the start of the proof of
Lemma~\ref{lem:theta}, it lies on the $\oarc J_1$-side of $J_2\cup L$ (so misses $E_2$) and on the $\oarc L$-side of $J_1\cup J_2$ (so misses $D$). Symmetrically
$E_2$ misses $D$. A point off the $\theta$-curve in none of the three sets would lie on the $\oarc L$-side of $J_1\cup J_2$, the $\oarc J_2$-side of $J_1\cup L$ and the
$\oarc J_1$-side of $J_2\cup L$, which Lemma~\ref{lem:theta}(a) excludes.

\emph{$U$ lies in $E_1$ or in $E_2$.} $U$ is connected and misses $M$, which contains the $\theta$-curve. It misses $D$, so lies on the side of $J_1\cup J_2$ containing $\oarc L$. By the $\theta$-curve lemma
(Lemma~\ref{lem:theta}(a)) it is not also on the $\oarc J_2$-side of $J_1\cup L$ and the $\oarc J_1$-side of $J_2\cup L$. Say $U\subseteq E_1$ (the other case is symmetric).

Put $X_1=M\cap(E_1\cup\oarc J_1)$ and $X_2=M\cap(E_2\cup\oarc J_2\cup\oarc L)$. As $M\cap D=\emptyset$, these sets cover $M\setminus\set{A,B}$; they are disjoint, and nonempty since $\oarc J_1\subseteq X_1$, $\oarc J_2\subseteq X_2$.
Since $\cl{E_2}\subseteq E_2\cup J_2\cup L$, the closure of $X_2$ misses $X_1$. Since $\cl{E_1}\subseteq E_1\cup J_1\cup L$, the closure of $X_1$ misses $E_2\cup\oarc J_2$; it remains to see that it misses $\oarc L$. This is the
difficulty: $\oarc L$ lies in the closure of $E_1$.

As $U\subseteq E_1$, $L'\subseteq\cl U\subseteq E_1\cup J_1\cup L$, so $L'$ meets $J_2$ only in $A$, $B$, and $L'$, $L$, $J_2$ form a second $\theta$-curve. Let $m\in M\cap E_1$, $m\notin L'$. Then $m$ is off this
$\theta$-curve; $m$ and $\oarc J_2$ lie outside $U$, so on the same side of $L\cup L'$; and $m$ and $\oarc L'$ lie outside $E_2$, so on the same side of $J_2\cup L$. By Lemma~\ref{lem:theta}(a), $m$ is not on the side
$S$ of $L'\cup J_2$ that contains $\oarc L$. Hence $M\cap E_1\setminus L'$ lies in the other side $S'$, whose closure lies in $S'\cup L'\cup J_2$ and misses $\oarc L$. And $M\cap E_1\cap L'$ lies in the closed set $L'$, which
misses $\oarc L$. So $\oarc L$ is not in the closure of $M\cap E_1$, and not in that of $M\cap\oarc J_1\subseteq J_1$.
\end{proof}

\formal{\lean{HPC.Jones.theta_separation}, \lean{HPC.Jones.theta_separation_core}, with
\lean{HPC.Jones.mem_cc_of_not_mem_domain} and \lean{HPC.Jones.side_eq_cc}.}

\begin{proposition}[Case 2]\label{prop:case2}
If $M\setminus J$ is not connected, some point of $M$ is a local separating point.
\end{proposition}

\begin{proof}
If $M=J$, every point of $M$ is a local separating point (Lemma~\ref{lem:local-sep}(a)). Otherwise $M\setminus J=H\cup K$ with $H$, $K$ nonempty and separated.

\emph{Step 1: an arc $T$.} $M\setminus J$ misses $D$, so it lies in the other component $E$ of $\R^2\setminus J$, which is open and connected. Join a point of $H$ to a point of $K$ by an arc in $E$ (a polygonal arc: connected
open subsets of the plane are polygonally connected), and trim it: from its last point $x$ on $\cl H$ to the next point $y$ on $\cl K$. Then $x\in H$, $y\in K$ (as $x,y\in E$ miss $J$), and the open arc $\oarc T$ between them misses $\cl H\cup\cl K\supseteq
M\setminus J$ and $J$, so it lies in a complementary domain $U$ of $M$, and $x,y\in\fr U$.

\emph{Step 2: the curve $C$.} $C=\fr U$ is a simple closed curve (Lemma~\ref{lem:jones49-setting}) through $x$ and $y$, and $U\ne D$ because $x\in\fr U\setminus J$.

\emph{Step 3.} Each of the two arcs of $C$ from $x$ to $y$ meets $J$: otherwise it would be a connected subset of $M\setminus J$ meeting $H$ and $K$.

\emph{Step 4: the arc $AxB$.} Let $A$ and $B$ be the first points of $J$ on the two arcs of $C$ from $x$ to $y$, starting at $x$. Then $A\ne B$, and $C=L\cup L'$ where $L$ (the arc $AxB$) and $L'$ are arcs from $A$ to $B$
with $L\cap L'=\set{A,B}$ and $L\cap J=\set{A,B}$.

\emph{Step 5.} $A$ and $B$ split $J$ into arcs $J_1$, $J_2$, and $J_1$, $J_2$, $L$ form a $\theta$-curve in $M$. By Lemma~\ref{lem:theta-sep}, $M\setminus\set{A,B}$ is the union of two nonempty separated sets.

\emph{Steps 6--7.} $M\setminus\set B$ is connected (Lemma~\ref{lem:jones49-setting}), so $A$ is a local separating point by Lemma~\ref{lem:local-sep}(b).
\end{proof}

\formal{\lean{HPC.Jones.exists_domain_meeting}, \lean{HPC.Jones.exists_theta_arcs},
\lean{HPC.Jones.jones49_case2_core}, \lean{HPC.Jones.jones49_case2}; the polygonal arc is
\lean{HPC.Jones.exists_tarc_of_isOpen_plane}. Jones: ``It follows that $M-(A+B)$ is \emph{not} connected, but $M-B$
\emph{is} connected. Hence $A$ is a local separating point of $M$.''}

\section{The theorems}

\begin{theorem}[Jones 1949, Theorem~1]\label{thm:jones49-1}
A nondegenerate homogeneous aposyndetic plane continuum is a simple closed curve.
\end{theorem}

\begin{proof}
If some point of $M$ has order $2$, $M$ is a simple closed curve by Corollary~\ref{cor:order-two}. Otherwise, with $x$, $D$, $J$ as above, Proposition~\ref{prop:case1} or
Proposition~\ref{prop:case2} gives a local separating point, and Corollary~\ref{cor:some-order-two} a point of order $2$, a contradiction.
\end{proof}

\formal{\lean{HPC.Jones.jones49_theorem1}, of type \lean{HPC.Jones.Jones49Thm1}.}

\begin{corollary}[Jones 1949, Theorem~2]\label{cor:jones49-2}
A nondegenerate homogeneous plane continuum without weak cut points is a simple closed curve.
\end{corollary}

\begin{proof}
By Jones's Theorem~18 of 1948 (Remark~\ref{rem:jones48-rest}), a continuum without weak cut points is aposyndetic at the points of a dense set, in particular at some point.
By homogeneity (Lemma~\ref{lem:apos-basic}(b)) it is aposyndetic at every point. Apply Theorem~\ref{thm:jones49-1}.
\end{proof}

\formal{\lean{HPC.Jones.jones49_theorem2}, from \lean{HPC.Jones.theorem18}; the one-point space is again a
counterexample to the printed statement, \lean{HPC.Jones.not_jones49_theorem2_point}.}

\section*{Exercises}

\begin{exercise}
Show that every hypothesis of Theorem~\ref{thm:jones49-1} is needed: give a homogeneous plane continuum that is not aposyndetic, an aposyndetic plane continuum that is not
homogeneous, and a homogeneous aposyndetic continuum (not in the plane) that is not a simple closed curve.
\end{exercise}

\begin{exercise}
Let $M$ be a $\theta$-curve with vertices $A$, $B$. Show that $M\setminus\set{A,B}$ is not connected, that $A$ is a local separating point and has order $3$,
and explain why this does not contradict Corollary~\ref{cor:some-order-two}.
\end{exercise}

\begin{exercise}
Where in the proof of Theorem~\ref{thm:jones49-1} is it used that $M$ lies in the plane? Find all the places.
\end{exercise}

%======================================================================
\chapter{The decomposable case}\label{ch:decomposable}
%======================================================================

This chapter assembles the results of the volume. Let $M$ be a decomposable homogeneous plane continuum, $G$ Jones's
decomposition of $M$ into the sets $\Ls x$ (Chapter~\ref{ch:jones55}), and $\pi\colon M\to M/G$ the quotient map. The plan is
Jones's (1955, Theorem~2 and the closing remarks), completed with the indecomposable case of Volume~6:
\begin{enumerate}
\item $M/G$ is homeomorphic to a plane continuum (Moore's Theorem~22);
\item hence $M/G$ is a simple closed curve (Jones 1949, Chapter~\ref{ch:jones49});
\item $M$ has exactly two complementary domains, and the elements are indecomposable;
\item the nondegenerate elements are pseudo-arcs (Volume~6);
\item so $M$ is a simple closed curve, or $M$ maps onto the circle by an open, atomic map whose fibres are pseudo-arcs.
\end{enumerate}

Throughout, $M\subseteq\R^2$ is a decomposable homogeneous continuum.

\section{The decomposition space lies in a plane}

\begin{definition}\label{def:gprime}
$G'$ is the decomposition of $\R^2$ whose elements are the sets $\Ls x$, $x\in M$, and the single points of $\R^2\setminus M$.
\end{definition}

\formal{\lean{HPC.Jones.gPrime}.}

\begin{proposition}[Jones 1955, p.~739]\label{prop:gprime}
$G'$ is a Moore decomposition of the plane (\vfive{Definition~2.1}). Consequently the space of elements of $G'$ is homeomorphic to $\R^2$,
and $M/G$ is homeomorphic to a compact connected subset of $\R^2$.
\end{proposition}

\begin{proof}
Every element is a continuum (Theorem~\ref{thm:jones48-3}) and none separates the plane (Theorem~\ref{thm:elt-not-separates}; points do not separate
the plane). For upper semicontinuity it suffices that the quotient map $\R^2\to G'$ is closed: the saturation of a closed set $F$ is $F$ together with
the saturation of the compact set $F\cap M$ under $G$, which is closed because $G$ is upper semicontinuous (Theorem~\ref{thm:jones55-usc},
\vthree{Theorem~6.5}(a)); and a closed quotient map with compact fibres gives Moore's distance form of upper semicontinuity by a compactness argument.
By Moore's Theorem~22 (\vfive{Theorem~9.13}) the space of elements of $G'$ is a plane. The inclusion $M\to\R^2$ induces a continuous injection of the
compact space $M/G$ into it, a homeomorphism onto its image, which is compact and connected.
\end{proof}

\formal{\lean{HPC.Jones.isClosedMap_mk_gPrime}, \lean{HPC.Jones.mooreDecomp_gPrime}, \lean{HPC.Jones.jones55_plane},
\lean{HPC.Jones.jones55_embedding_data}, \lean{HPC.Jones.jones55_embedding}. This is the construction of \vfive{Corollary~9.16}; Jones
states the upper semicontinuity of $G'$ without argument.}

\section{Jones's Theorem 2}

\begin{theorem}[Jones 1955, Theorem~2]\label{thm:jones55-2}
Let $M\subseteq\R^2$ be a decomposable homogeneous continuum. Then $G$ is a continuous decomposition of $M$, and
\begin{enumerate}[label=(\arabic*)]
\item $M/G$ is a simple closed curve;
\item each element of $G$ is a homogeneous continuum that does not separate the plane.
\end{enumerate}
\end{theorem}

\begin{proof}
Continuity is Theorem~\ref{thm:jones55-usc}. (1) By Theorem~\ref{thm:jones55-a}, $M/G$ is a nondegenerate homogeneous aposyndetic continuum, and by
Proposition~\ref{prop:gprime} it is homeomorphic to a plane continuum $K$. Homogeneity and aposyndesis are topological properties, so $K$ is a nondegenerate
homogeneous aposyndetic plane continuum, hence a simple closed curve by Theorem~\ref{thm:jones49-1}. (2) Proposition~\ref{prop:jones55-elements} and
Theorem~\ref{thm:elt-not-separates}.
\end{proof}

\formal{\lean{HPC.Jones.jones55_circle}, \lean{HPC.Jones.jones55_theorem2'} (and \lean{HPC.Jones.jones55_theorem2}, with
Theorem~\ref{thm:jones49-1} as a hypothesis).}

\begin{corollary}[Jones 1955, closing remarks]\label{cor:jones55-remarks}
\leavevmode
\begin{enumerate}[label=(\alph*)]
\item A decomposable homogeneous plane continuum has exactly two complementary domains.
\item The elements of $G$ are indecomposable.
\item A homogeneous plane continuum with more than two complementary domains is indecomposable.
\end{enumerate}
\end{corollary}

\begin{proof}
(a) In the plane of elements of $G'$ (Proposition~\ref{prop:gprime}) the image of $M$ is a simple closed curve (Theorem~\ref{thm:jones55-2}(1)). By the
Jordan curve theorem (\vfive{Theorem~4.8}) its complement consists of two disjoint connected open sets of elements. The union of the elements in each
of them is open (the quotient map is continuous), connected (\vfive{Proposition~2.5}(a)), and these two sets are disjoint and have union $\R^2\setminus M$, because
the elements of $G'$ outside $M$ are the points of $\R^2\setminus M$.

(b) If an element $\Ls x$ were decomposable, it would be a decomposable homogeneous plane continuum (Proposition~\ref{prop:jones55-elements}), so by (a) its complement
would be the union of two disjoint nonempty open sets: $\Ls x$ would separate the plane, against Theorem~\ref{thm:jones55-2}(2).

(c) is the contrapositive of (a): each complementary domain of $M$ is one of the two sets of (a).
\end{proof}

\formal{\lean{HPC.Jones.jones55_two_domains'}, \lean{HPC.Jones.jones55_elements_indecomposable'},
\lean{HPC.Jones.jones55_three_domains'}; the plane Jordan curve theorem used is \lean{HPC.Jones.jordan_curve_theorem}. Jones
derives (b) from his 1951 paper; the formal proof derives it from (a) as above.}

\section{The elements are pseudo-arcs}

\begin{proposition}\label{prop:elements-pseudoarc}
Every nondegenerate element $\Ls x$ of $G$ is a pseudo-arc (\vone{Definition~2.14}).
\end{proposition}

\begin{proof}
$\Ls x$ is homeomorphic to the plane continuum $\Ls x\subseteq\R^2$, which is homogeneous (Theorem~\ref{thm:jones55-2}(2)), indecomposable
(Corollary~\ref{cor:jones55-remarks}(b)) and nondegenerate. By \vsix{Theorem~7.4} (proved unconditionally in \vsix{Corollary~10.6}), it is a pseudo-arc.
\end{proof}

\formal{\lean{HPC.BJ59.isPseudoArc_Lset'}, from \lean{HPC.HO16.isPseudoArc_of_indec_homogeneous} and
\lean{HPC.HO16.ho16_theorem20}.}

\section{The main theorem}

\begin{definition}\label{def:atomic}
Let $f\colon X\to Y$ be a map of compact metric spaces.
\begin{enumerate}[label=(\alph*)]
\item $f$ is \emph{atomic} if each fibre $f^{-1}(y)$ is terminal: every continuum in $X$ meeting $f^{-1}(y)$ lies in it or contains it.
\item An \emph{atomic circle-of-pseudo-arcs map} is a continuous, open, surjective, atomic map $f\colon X\to\Sph$ whose fibres are pseudo-arcs.
\end{enumerate}
\end{definition}

\formal{\lean{HPC.BJ59.IsAtomicMap}, \lean{HPC.BJ59.IsAtomicCircPAMap}. The terminology is that of Domokos and Prajs.}

\begin{lemma}\label{lem:open-quotient}
The quotient map of a continuous decomposition of a compact metric space is open.
\end{lemma}

\begin{proof}
Let $U$ be open. We show that the set $U'$ of points whose element meets $U$ is open. Let $z\in U'$ and $y\in\elt z\cap U$. By lower semicontinuity at $\elt z$
there is an open $W\supseteq\elt z$ such that every element meeting $W$ meets $U$; so $W\subseteq U'$. As $U'$ is saturated, $\pi(U')=\pi(U)$ is open in the quotient.
\end{proof}

\formal{\lean{HPC.BJ59.isOpenMap_quotient_mk_of_continuous}.}

\begin{theorem}[The decomposable case]\label{thm:decomposable}
Let $M$ be a decomposable homogeneous plane continuum. Then either $M$ is homeomorphic to the circle, or there is an atomic circle-of-pseudo-arcs map
$f\colon M\to\Sph$, whose fibres are the elements $\Ls x$ of Jones's decomposition.
\end{theorem}

\begin{proof}
Let $e\colon M/G\to\Sph$ be a homeomorphism (Theorem~\ref{thm:jones55-2}(1)) and $q=e\circ\pi$. Then $q$ is continuous and surjective, and its fibres are the elements:
$q^{-1}(q(x))=\Ls x$.

If $M$ is aposyndetic, every element is a point (Proposition~\ref{prop:jones55-elements}(a)), so $q$ is a continuous bijection from a compact space onto a Hausdorff space, hence
a homeomorphism.

Otherwise every element is nondegenerate (Proposition~\ref{prop:jones55-elements}(b)), so every fibre of $q$ is a pseudo-arc (Proposition~\ref{prop:elements-pseudoarc}). $q$ is open
by Lemma~\ref{lem:open-quotient}, since $G$ is continuous (Theorem~\ref{thm:jones55-usc}). And $q$ is atomic: a continuum $L$ meeting $\Ls x$ either lies in $\Ls x$ or meets its
complement, and then contains $\Ls x$ by Theorem~\ref{thm:jones55-c}.
\end{proof}

\formal{\lean{HPC.BJ59.circle_or_atomic_of_decomposable}.}

\begin{remark}[The circle-like form]\label{rem:circle-like}
A \emph{circle of pseudo-arcs} in the sense of Bing and Jones (1959) is a circle-like continuum with a continuous open map onto $\Sph$ whose fibres are pseudo-arcs. The proof
of Theorem~\ref{thm:decomposable} gives such a map whenever $M$ is known to be circle-like (\vthree{Chapter~5}). That a decomposable homogeneous plane continuum which is not a simple
closed curve is circle-like, and that any two circles of pseudo-arcs are homeomorphic, are the subject of Volume~8; there the atomic map of Theorem~\ref{thm:decomposable} is the
starting point.
\end{remark}

\formal{\lean{HPC.BJ59.circle_or_circPAMap_of_decomposable}, with the map structure \lean{HPC.BJ59.IsCircPAMap}.}

\section*{Exercises}

\begin{exercise}
Show that the projection $P\times\Sph\to\Sph$, $P$ a pseudo-arc, is open and has pseudo-arc fibres, but is not atomic. Using that $P\times\Sph$ is homogeneous
(\vtwo{Theorem~8.4}) and aposyndetic (Theorem~\ref{thm:jones48-7}), and Theorem~\ref{thm:decomposable}, show that $P\times\Sph$ does not embed in the plane.
\end{exercise}

\begin{exercise}
Let $f\colon X\to\Sph$ be an atomic circle-of-pseudo-arcs map. Show that the fibres of $f$ are exactly the sets $\Ls x$ of $X$. (Use that
the preimage of a closed arc of $\Sph$ is a continuum, and that the fibres have empty interior because $f$ is open.)
\end{exercise}

\begin{exercise}
Where does the proof of Theorem~\ref{thm:decomposable} use the plane? List the results of this volume that are used and need it.
\end{exercise}


\appendix
%======================================================================
\chapter{Correspondence with the formal proof}\label{app:formal}
%======================================================================

The table lists the definitions and results of this volume that carry a ``Formal counterpart''
line, together with the corresponding declarations of the Lean development. All declaration
names start with \lean{HPC.}. The formal statements are sometimes slightly more general, or are
split differently into lemmas; the table points to the place where the content of each result is
proved. Rows are in the order of the text; rows labelled ``Text of \dots'' belong to formal counterpart lines in the
running text of the named section.

The main theorems of the volume are, in the Lean development (prefix \lean{HPC.} omitted):
\begin{center}
\begin{tabular}{@{}p{0.42\textwidth}p{0.54\textwidth}@{}}
\toprule
Jones 1948, Theorem~3 (Theorem~\ref{thm:jones48-3}) & \lean{Jones.theorem3}\\
Jones 1948, Theorem~8 (Theorem~\ref{thm:jones48-8}) & \lean{Jones.theorem8}\\
Spreading (Theorem~\ref{thm:spreading}) & \lean{Jones.aposAt_of_isContinuum_interior}\\
Jones 1951, Lemma (Theorem~\ref{thm:jones51-lemma}) & \lean{Jones.jones51_lemma}\\
Bing 1948, Theorem~5 (Theorem~\ref{thm:bing5}) & \lean{Plane.arcs_theorem5_point}\\
Jones 1955, Theorem~1(b), (c) (Theorems~\ref{thm:jones55-b}, \ref{thm:jones55-c}) & \lean{Jones.jones55_b}, \lean{Jones.jones55_c}\\
$G$ is continuous (Theorem~\ref{thm:jones55-usc}) & \lean{Jones.jones55_continuous}\\
Jones 1955, Theorem~1(a) (Theorem~\ref{thm:jones55-a}) & \lean{Jones.jones55_a}\\
Elements do not separate (Theorem~\ref{thm:elt-not-separates}) & \lean{Jones.jones55_elt_not_separates}\\
Jones 1951, Theorem~1 (Theorem~\ref{thm:jones51-1}) & \lean{Jones.jones51_theorem1}\\
Menger's theorem (Theorem~\ref{thm:menger}) & \lean{Jones.menger_order_two}\\
Separating point--order theorem (Theorem~\ref{thm:sep-order}) & \lean{Jones.countable_sepPoint_not_hasOrder_two}\\
Whyburn (9.2) (Corollary~\ref{cor:whyburn}(a)) & \lean{Jones.whyburn_9_2}\\
Sector theorem (Theorem~\ref{thm:sector}) & \lean{Jones.sector_theorem}\\
Cyclic order (Theorem~\ref{thm:cyclic}) & \lean{Jones.sector_cyclic_order}\\
Whyburn 1939, Theorem~14 (Theorem~\ref{thm:whyburn14}) & \lean{Jones.whyburn39_theorem14}\\
Whyburn 1939, Corollary~1 (Theorem~\ref{thm:whyburn-cor1}) & \lean{Jones.whyburn39_corollary1}\\
Jones 1949, Theorem~1 (Theorem~\ref{thm:jones49-1}) & \lean{Jones.jones49_theorem1}\\
Jones 1949, Theorem~2 (Corollary~\ref{cor:jones49-2}) & \lean{Jones.jones49_theorem2}\\
Jones 1955, Theorem~2 (Theorem~\ref{thm:jones55-2}) & \lean{Jones.jones55_theorem2'}\\
Two complementary domains (Corollary~\ref{cor:jones55-remarks}) & \lean{Jones.jones55_two_domains'}, \lean{Jones.jones55_elements_indecomposable'}\\
The elements are pseudo-arcs (Proposition~\ref{prop:elements-pseudoarc}) & \lean{BJ59.isPseudoArc_Lset'}\\
The decomposable case (Theorem~\ref{thm:decomposable}) & \lean{BJ59.circle_or_atomic_of_decomposable}\\
\bottomrule
\end{tabular}
\end{center}
All of them are proved in the Lean development without \lean{sorry} and without axioms beyond the
standard ones of Lean and Mathlib. The printed forms that are false (Section~\ref{sec:intro-diff}) are
refuted by the declarations whose names begin with \lean{not_}.

\begin{longtable}{@{}p{0.40\textwidth}p{0.56\textwidth}@{}}
\toprule
In these notes & In the Lean development\\
\midrule
\endhead
\bottomrule
\endfoot
\multicolumn{2}{@{}l}{\emph{Chapter~\ref{ch:apos}}}\\
Text of ``Conventions and definitions'' & The setting is a type \lean{X} with \lean{[MetricSpace X] [CompactSpace X] [ConnectedSpace X]}; continua in $X$ are \lean{HPC.Fold.Q78.IsContinuum}; decomposability is \lean{HPC.Jones.IsDecomposableSpace}.\\
Definition~\ref{def:apos} & \lean{HPC.Jones.AposAt}, \lean{HPC.Jones.AposyndeticAt}, \lean{HPC.Jones.IsAposyndetic}, \lean{HPC.Jones.TotallyNonAposyndetic}; $\Us x$ is \lean{HPC.Jones.aposSet} (with \lean{HPC.Jones.aposSet_eq}) and $\Ls x$ is \lean{HPC.Jones.Lset} (with \lean{HPC.Jones.Lset_compl}).\\
Definition~\ref{def:slc} & \lean{HPC.Jones.SemiLocConnAt}.\\
Definition~\ref{def:cuts} & \lean{HPC.Jones.CutsPts}, \lean{HPC.Jones.IsWeakCutPoint}, \lean{HPC.Jones.IsStrongCutPoint}.\\
Lemma~\ref{lem:apos-basic} & \lean{HPC.Jones.aposAt_of_mem}, \lean{HPC.Jones.aposAt_homeomorph_iff}, \lean{HPC.Jones.aposSet_homeomorph}, \lean{HPC.Jones.Lset_homeomorph}.\\
Theorem~\ref{thm:jones48-1} & \lean{HPC.Jones.isOpen_setOf_aposAt}, \lean{HPC.Jones.isOpen_aposSet}, \lean{HPC.Jones.isClosed_Lset}, \lean{HPC.Jones.theorem1}. Jones states the sequential form with the hypotheses $x\ne y$, $x_i\ne y_i$, which are not needed.\\
Lemma~\ref{lem:heine-borel} & \lean{HPC.Jones.exists_finite_apos_cover}, \lean{HPC.Jones.exists_nhd_uniform_apos}.\\
Lemma~\ref{lem:apos-frontier} & \lean{HPC.Jones.aposAt_of_frontier}; the last step is \lean{HPC.Jones.exists_open_subset_component}.\\
Theorem~\ref{thm:jones48-3} & \lean{HPC.Jones.theorem3}, \lean{HPC.Jones.isContinuum_Lset}. Jones prints no proof.\\
Theorem~\ref{thm:jones48-0} & \lean{HPC.Jones.theorem0}. Jones only says that it ``may be established with the help of the Heine--Borel theorem''.\\
Corollary~\ref{cor:jones41} & \lean{HPC.Jones.semiLocConn_of_isAposyndetic}, \lean{HPC.Jones.isAposyndetic_of_semiLocConn}.\\
Proposition~\ref{prop:whyburn62} & \lean{HPC.Jones.whyburn_6_2}; also Whyburn's (6.1) and (6.21), \lean{HPC.Jones.whyburn_6_1}, \lean{HPC.Jones.whyburn_6_21}. Whyburn argues with the components of $X\setminus V_\delta(x)$; the formal proof goes through Theorem~\ref{thm:jones48-0}.\\
Theorem~\ref{thm:jones48-8} & \lean{HPC.Jones.theorem8}, \lean{HPC.Jones.theorem9}; the first direction is \lean{HPC.Jones.aposAt_of_decomposition}.\\
Theorem~\ref{thm:jones48-7} & \lean{HPC.Jones.theorem7}, with \lean{HPC.Jones.isContinuum_prod_cross}.\\
Remark~\ref{rem:jones48-rest} & \lean{HPC.Jones.theorem6}, \lean{HPC.Jones.not_theorem6_printed}, \lean{HPC.Jones.theorem10}, \lean{HPC.Jones.theorem10_printed_iff}, \lean{HPC.Jones.not_theorem10_printed_unitInterval}, \lean{HPC.Jones.theorem11}, \lean{HPC.Jones.theorem12}, \lean{HPC.Jones.theorem13}, \lean{HPC.Jones.totallyNonApos_exists_weakCutPoint}, \lean{HPC.Jones.theorem18} (in \lean{HPC/Jones/Basic.lean}, \lean{HPC/Jones/Theorem6.lean}, \lean{HPC/Jones/Theorem13.lean}, \lean{HPC/Jones/CutPoints.lean}).\\[4pt]
\multicolumn{2}{@{}l}{\emph{Chapter~\ref{ch:spread}}}\\
Text of ``chapter introduction'' & \lean{HPC.Jones.IsHomogeneous}; it is the hypothesis of \lean{HPC.Effros.effros_homogeneous} (\lean{HPC.Jones.isHomogeneous_iff_isPretransitive}).\\
Theorem~\ref{thm:spreading} & \lean{HPC.Jones.aposAt_of_isContinuum_interior}, from \lean{HPC.Effros.effros_homogeneous}.\\
Corollary~\ref{cor:trans} & \lean{HPC.Jones.not_aposAt_trans}.\\
Lemma~\ref{lem:open-subset-U} & \lean{HPC.Jones.isOpen_setOf_subset_aposSet}.\\
Theorem~\ref{thm:jones51-lemma} & \lean{HPC.Jones.jones51_lemma}, with the function $g$ from \lean{HPC.Jones.exists_lsc_strictMono_aposSet}. Jones's proof runs a transfinite sequence of points $x_0,x_1,\dots$ through the countable ordinals with $\Us{x_i}$ strictly decreasing, and contradicts second countability; the lower semicontinuous function replaces the transfinite recursion.\\
Lemma~\ref{lem:tight-conn} & \lean{HPC.Plane.Arcs48.isPreconnected_closure_diff_union} (Bing: ``$F_i+N$ is a continuum'').\\
Lemma~\ref{lem:tight-piece} & \lean{HPC.Plane.Arcs48.exists_nowhereDense_piece}.\\
Lemma~\ref{lem:tight-case} & \lean{HPC.Plane.Arcs48.exists_point_of_tight}.\\
Lemma~\ref{lem:bing-seq} & \lean{HPC.Plane.Arcs48.tight_or_point}, \lean{HPC.Plane.Arcs48.exists_point_dense}.\\
Theorem~\ref{thm:bing5} & \lean{HPC.Plane.arcs_theorem5_point}; the general form for a proper subset $R$ in place of $\set r$ is \lean{HPC.Plane.arcs_theorem5}.\\
Corollary~\ref{cor:bing2} & \lean{HPC.Plane.arcs_corollary2}, with Bing's relation ``$R$ cuts $A$ from $B$ in $M$'' as \lean{HPC.Plane.Cuts}.\\
Proposition~\ref{prop:claim35} & \lean{HPC.Jones.exists_continuum_meets_of_open}, \lean{HPC.Jones.jones55_claim3}, \lean{HPC.Jones.exists_continuum_of_mem_aposSet}, \lean{HPC.Jones.jones55_claim5}.\\[4pt]
\multicolumn{2}{@{}l}{\emph{Chapter~\ref{ch:jones55}}}\\
Lemma~\ref{lem:U-eq} & \lean{HPC.Jones.aposSet_eq_of_not_aposAt}.\\
Theorem~\ref{thm:jones55-b} & \lean{HPC.Jones.mem_Lset_iff_Lset_eq}, \lean{HPC.Jones.jones55_b}. Jones assumes $X$ decomposable, which is not needed.\\
Theorem~\ref{thm:jones55-c} & \lean{HPC.Jones.Lset_subset_of_meets}, \lean{HPC.Jones.jones55_c}.\\
Text of ``The sets $\Ls x$ partition $X$'' & The relation $\Ls x=\Ls y$ is \lean{HPC.Jones.lRel}, the setoid is \lean{HPC.Jones.lSetoid}, and \lean{HPC.Jones.cls_lSetoid} says that its classes are the sets $\Ls x$.\\
Remark~\ref{rem:props} & \lean{HPC.Jones.jones55_prop_a}, \lean{HPC.Jones.jones55_prop_b}, \lean{HPC.Jones.jones55_prop_c}, \lean{HPC.Jones.jones55_prop_d}, \lean{HPC.Jones.jones55_prop_e}, \lean{HPC.Jones.jones55_prop_f}, with $N_o$ as \lean{HPC.Jones.Nset} and \lean{HPC.Jones.Nset_eq_Lset}. Jones's intermediate claims (1)--(12) are also proved, as \lean{HPC.Jones.jones55_claim1} through \lean{HPC.Jones.jones55_claim12}.\\
Theorem~\ref{thm:jones55-usc} & \lean{HPC.Jones.isClosed_lSetoid}, \lean{HPC.Jones.jones55_usc}, \lean{HPC.Jones.lSetoid_homeomorph}, \lean{HPC.Jones.jones55_continuous}; from \lean{HPC.Jones.isUSCDecomp_of_isClosed}, \lean{HPC.Jones.exists_lscAt} (with \lean{HPC.Jones.residual_lscAt}) and \lean{HPC.Jones.lscAt_homeomorph}. Jones argues upper semicontinuity with convergent sequences of elements and property (f).\\
Theorem~\ref{thm:jones55-a} & \lean{HPC.Jones.jones55_a}; homogeneity from \lean{HPC.Jones.exists_homeomorph_quotient}, nondegeneracy from \lean{HPC.Jones.theorem8}. Jones states aposyndesis of $X/G$ without argument; the proof above uses the continuity of $G$.\\
Proposition~\ref{prop:jones55-elements} & \lean{HPC.Jones.Lset_eq_singleton_of_isAposyndetic}, \lean{HPC.Jones.jones55_elements}, \lean{HPC.Jones.isHomogeneous_Lset}, \lean{HPC.Jones.aposSet_nonempty}, \lean{HPC.Jones.interior_Lset_eq_empty}, \lean{HPC.Jones.jones55_claim12}.\\
Proposition~\ref{prop:compl-connected} & \lean{HPC.Jones.jones55_compl_continuumwise}, \lean{HPC.Jones.jones55_compl_connected}; Jones writes ``it is clear that for each $g$ of $G$, $M-g$ is connected (strongly or continuum-wise connected) because of (5)''.\\
Lemma~\ref{lem:disjoint-open-countable} & \lean{HPC.Jones.countable_of_pairwiseDisjoint_open}.\\
Theorem~\ref{thm:elt-not-separates} & \lean{HPC.Jones.jones55_elt_not_separates}, using \lean{HPC.Jones.separatesPlane_of_homeomorph}, \lean{HPC.Jones.closure_component_compl_subset}, \lean{HPC.Jones.uncountable_quotient} and \lean{HPC.Jones.uncountable_of_connected_metrizable}. Jones's proof says only that ``if one separates, all do, which is impossible''; the disjointness of the domains $W_g$ is the missing step.\\
Theorem~\ref{thm:jones51-1} & \lean{HPC.Jones.jones51_theorem1}. Jones's own proof builds the decomposition of the 1955 paper in an early form, passes to the decomposition space and applies Whyburn's (6.21) there.\\[4pt]
\multicolumn{2}{@{}l}{\emph{Chapter~\ref{ch:order}}}\\
Definition~\ref{def:order} & \lean{HPC.Jones.OrderLE}, \lean{HPC.Jones.HasOrder}.\\
Lemma~\ref{lem:order-basic} & \lean{HPC.Jones.orderLE_homeomorph}, \lean{HPC.Jones.hasOrder_homeomorph}, \lean{HPC.Jones.orderLE_subtype}.\\
Lemma~\ref{lem:triod} & \lean{HPC.Jones.not_orderLE_two_of_triod}.\\
Proposition~\ref{prop:finite-order-lc} & \lean{HPC.Jones.locallyConnectedSpace_of_orderLE}.\\
Lemma~\ref{lem:open-continuum} & \lean{HPC.Jones.exists_continuum_of_isOpen}.\\
Proposition~\ref{prop:finite-order-arcs} & \lean{HPC.Jones.exists_tarc_of_orderLE}.\\
Lemma~\ref{lem:legs} & \lean{HPC.Jones.arc_eq_of_legs}.\\
Theorem~\ref{thm:menger} & \lean{HPC.Jones.menger_order_two}, with \lean{HPC.Jones.frontier_connectedComponentIn_subset}, \lean{HPC.Jones.isPreconnected_diff_of_frontier} and \lean{HPC.Jones.tscc_homeomorph_circle} (two arcs glued along their ends, \lean{HPC.Jones.glue_circle}). Menger's own route goes through his $n$-Beinsatz and the classification of curves of order~2; the proof above is his second argument (the first point on an arc is a branch point), applied twice.\\
Corollary~\ref{cor:order-two} & \lean{HPC.Jones.jones49_order_two}.\\[4pt]
\multicolumn{2}{@{}l}{\emph{Chapter~\ref{ch:whyburn}}}\\
Definition~\ref{def:sep-point} & \lean{HPC.Jones.IsSepPointOf}, \lean{HPC.Jones.IsLocalSepPoint}. Whyburn's phrase is ``$R-p=R_a+R_b$ with $a\in R_a$, $b\in R_b$ in the component of $R$ containing $p$''.\\
Text of ``Separating points of an open set'' & \lean{HPC.Jones.IsLocalSepPoint.homeomorph}.\\
Lemma~\ref{lem:side} & \lean{HPC.Jones.isPreconnected_side}, \lean{HPC.Jones.isPreconnected_diff_side}.\\
Proposition~\ref{prop:whyburn91} & \lean{HPC.Jones.whyburn_9_1}, with \lean{HPC.Jones.exists_continuum_bump_compact} (boundary bumping for a compact piece of a connected set) and \lean{HPC.Jones.whyburn_9_1_side}. Whyburn uses at this point that ``the components \dots\ are nondegenerate'' (his I, (10.2)).\\
Lemma~\ref{lem:sep-not-order1} & \lean{HPC.Jones.not_orderLE_one_of_isSepPointOf}, \lean{HPC.Jones.hasOrder_two_of_isSepPointOf}.\\
Definition~\ref{def:sepchain} & \lean{HPC.Jones.SepChain}, with the order \lean{HPC.Jones.SepChain.lt}.\\
Lemma~\ref{lem:chain-order} & \lean{HPC.Jones.SepChain.lt_iff}, \lean{HPC.Jones.SepChain.lt_total}, \lean{HPC.Jones.SepChain.lt_asymm}, \lean{HPC.Jones.SepChain.lt_trans}, \lean{HPC.Jones.SepChain.B_sub}, \lean{HPC.Jones.SepChain.A_sub}.\\
Text of ``Separation chains'' & \lean{HPC.Jones.SepChain.M}, \lean{HPC.Jones.SepChain.M_isPreconnected}, \lean{HPC.Jones.SepChain.M_closed}, \lean{HPC.Jones.SepChain.M_mono_left}, \lean{HPC.Jones.SepChain.M_mono_right}, \lean{HPC.Jones.SepChain.M_sub_BA}.\\
Lemma~\ref{lem:counting} & \lean{HPC.Jones.countable_not_mem_closure_of_lt}.\\
Text of ``Separation chains'' & \lean{HPC.Jones.SepChain.Good}, \lean{HPC.Jones.SepChain.Res}.\\
Lemma~\ref{lem:good-res} & \lean{HPC.Jones.SepChain.countable_not_good}, \lean{HPC.Jones.SepChain.Res_sub_B}, \lean{HPC.Jones.SepChain.Res_sub_A}, \lean{HPC.Jones.SepChain.countable_res}.\\
Proposition~\ref{prop:good-order2} & \lean{HPC.Jones.SepChain.exists_small_interval}, \lean{HPC.Jones.SepChain.exists_open_frontier_pair} (with the cut-wire lemma \lean{HPC.Jones.exists_clopen_of_connectedComponentIn_disjoint}), \lean{HPC.Jones.SepChain.orderLE_two}. This replaces Whyburn's (8.3) and (2.21)--(2.22).\\
Corollary~\ref{cor:chain-countable} & \lean{HPC.Jones.SepChain.countable_of_not_orderLE_two}.\\
Theorem~\ref{thm:sep-order} & \lean{HPC.Jones.countable_sepPoint_not_hasOrder_two}, from \lean{HPC.Jones.countable_sepPoint_not_orderLE_two}, \lean{HPC.Jones.sepT}, \lean{HPC.Jones.sepT_mem_component}, \lean{HPC.Jones.countable_sepT}.\\
Corollary~\ref{cor:whyburn} & \lean{HPC.Jones.whyburn_9_2}, \lean{HPC.Jones.whyburn_8_4}, \lean{HPC.Jones.whyburn_3_2}.\\
Corollary~\ref{cor:some-order-two} & \lean{HPC.Jones.jones49_some_order_two}, \lean{HPC.Jones.jones49_some_order_two_of_whyburn_9_2}.\\[4pt]
\multicolumn{2}{@{}l}{\emph{Chapter~\ref{ch:sectors}}}\\
Definition~\ref{def:annulus} & \lean{HPC.Jones.Sector.Ann}, \lean{HPC.Jones.Sector.IsCrossA}, \lean{HPC.Jones.Sector.Strip}, \lean{HPC.Jones.Sector.IsCross}, \lean{HPC.Jones.Sector.Ecov}, \lean{HPC.Jones.Sector.dist_Ecov}.\\
Lemma~\ref{lem:lift} & \lean{HPC.Jones.Sector.exists_log_cross}, \lean{HPC.Jones.Sector.exists_lift} (the lift as \lean{HPC.Jones.Sector.IsLift}), \lean{HPC.Jones.Sector.exists_normalized_lift} (a lift whose leftmost points have first coordinate in $[0,2\pi)$).\\
Definition~\ref{def:sides} & \lean{HPC.Jones.Sector.Rc}, \lean{HPC.Jones.Sector.Lc}, \lean{HPC.Jones.Sector.NotR}, with \lean{HPC.Jones.Sector.SideOf}, \lean{HPC.Jones.Sector.FarR}, \lean{HPC.Jones.Sector.FarL}.\\
Text of ``The two sides of a crossing set'' & \lean{HPC.Jones.Sector.sideOf_absorb}, \lean{HPC.Jones.Sector.sideOf_open_rel}, \lean{HPC.Jones.Sector.GoodSide.exists_mem_closure}, \lean{HPC.Jones.Sector.GoodSide.union_isPreconnected}.\\
Theorem~\ref{thm:strip-sep} & \lean{HPC.Jones.Sector.no_cross_join}, \lean{HPC.Jones.Sector.exists_compact_join}, \lean{HPC.Jones.Sector.lc_disjoint_rc}.\\
Definition~\ref{def:prec} & \lean{HPC.Jones.Sector.Prec}.\\
Proposition~\ref{prop:prec} & \lean{HPC.Jones.Sector.prec_total}, \lean{HPC.Jones.Sector.prec_asymm}, \lean{HPC.Jones.Sector.prec_trans}, \lean{HPC.Jones.Sector.prec_of_notR}, \lean{HPC.Jones.Sector.rc_mono}, \lean{HPC.Jones.Sector.notR_mono}, \lean{HPC.Jones.Sector.prec_min}, \lean{HPC.Jones.Sector.prec_sh_self}.\\
Text of ``Sectors'' & \lean{HPC.Jones.Sector.LiftFam}, \lean{HPC.Jones.Sector.Sec}, \lean{HPC.Jones.Sector.sec_open_rel}, \lean{HPC.Jones.Sector.LiftFam.prec_shift_pos}.\\
Proposition~\ref{prop:sectors} & \lean{HPC.Jones.Sector.LiftFam.sub_sec}, \lean{HPC.Jones.Sector.LiftFam.sec_disjoint_mem}, \lean{HPC.Jones.Sector.LiftFam.sec_nonempty}, \lean{HPC.Jones.Sector.LiftFam.sec_disjoint}, \lean{HPC.Jones.Sector.LiftFam.sec_cover}.\\
Theorem~\ref{thm:sector} & \lean{HPC.Jones.sector_theorem}, from \lean{HPC.Jones.Sector.sector_two} and \lean{HPC.Jones.Sector.two_piece_lift}; Moore uses it on p.~236 of his 1922 paper (``must clearly contain a point either of $k_{n-1}$ or of $k_{n+1}$'').\\
Theorem~\ref{thm:cyclic} & \lean{HPC.Jones.sector_cyclic_order}, from \lean{HPC.Jones.Sector.sector_cyclic}, \lean{HPC.Jones.Sector.sector_props}, \lean{HPC.Jones.Sector.preconnected_sub_piece}. Moore (p.~237) orders the continua by their points on a middle circle; the order of the lifts avoids having to show that this is well defined.\\
Definition~\ref{def:cik} & \lean{HPC.Jones.CikAt}.\\
Proposition~\ref{prop:cik} & \lean{HPC.Jones.locallyConnectedSpace_of_cikAt}, \lean{HPC.Jones.exists_crossing_continua}, and for closed sets \lean{HPC.Jones.Sector.exists_crossing_continua_closed}. Moore attributes (b) to ``the argument beginning with the second line from the bottom of page 365'' of a 1918 note.\\[4pt]
\multicolumn{2}{@{}l}{\emph{Chapter~\ref{ch:boundaries}}}\\
Text of ``chapter introduction'' & \lean{HPC.Jones.isConnected_frontier_complDomain}.\\
Lemma~\ref{lem:whyburn622} & \lean{HPC.Jones.whyburn39_6_22_sec14}.\\
Lemma~\ref{lem:join} & \lean{HPC.Jones.Sec14.join_in_domain}; it replaces Whyburn's $\theta$-curve step and Moore's curve $C$ (\lean{HPC.Jones.whyburn39_moore_curve} records the form used).\\
Proposition~\ref{prop:local-join} & \lean{HPC.Jones.Sec14.local_join_of_noncut}, \lean{HPC.Jones.whyburn39_local_noncut}.\\
Text of ``Theorem 14'' & \lean{HPC.Jones.Sec14.BadAt}.\\
Lemma~\ref{lem:accum} & \lean{HPC.Jones.Sec14.not_orderLE_of_accumulating}, recorded as \lean{HPC.Jones.whyburn39_curve_from_order_two}; it replaces Whyburn's curve $K$ and his pigeonhole argument.\\
Lemma~\ref{lem:bad-point} & \lean{HPC.Jones.Sec14.exists_bad_point_of_crossing}. This is Moore's 1922 argument (p.~236) at an arbitrary radius; at the middle radius it gives Moore's local sufficiency theorem \lean{HPC.Jones.moore22_sufficiency_local}, \lean{HPC.Jones.Sector.sufficiency_local}.\\
Theorem~\ref{thm:whyburn14} & \lean{HPC.Jones.whyburn39_theorem14}.\\
Lemma~\ref{lem:theta} & \lean{HPC.Jones.theta_no_common_side}, \lean{HPC.Jones.theta_not_subset_closure} (with \lean{HPC.Jones.jordan_side}).\\
Lemma~\ref{lem:frontier-cyclic} & \lean{HPC.Jones.frontier_cyclic}, \lean{HPC.Jones.frontier_cyclic'}.\\
Proposition~\ref{prop:cyclic-peano} & \lean{HPC.Jones.exists_tarc_in_open}, \lean{HPC.Jones.exists_tarc_of_locallyConnected}, \lean{HPC.Jones.exists_scc_of_cyclic}, \lean{HPC.Jones.exists_theta_of_cyclic}, \lean{HPC.Jones.theta_of_cyclic_peano}.\\
Theorem~\ref{thm:whyburn-cor1} & \lean{HPC.Jones.whyburn39_corollary1}, with \lean{HPC.Jones.tscc_homeomorph_circle}. Whyburn gives no proof and refers to Wilder (1938). The hypothesis ``no cut points'' is stated there as: $M\setminus\set z$ is connected and nonempty for every $z\in M$, which includes nondegeneracy. The same argument also shows that every point of $B$ is accessible from $R$: \lean{HPC.Jones.accessible_of_cyclic}, which the classification does not need.\\[4pt]
\multicolumn{2}{@{}l}{\emph{Chapter~\ref{ch:jones49}}}\\
Text of ``chapter introduction'' & The statement is \lean{HPC.Jones.Jones49Thm1}; the one-point counterexample to the printed version is \lean{HPC.Jones.not_jones49_theorem1_point}.\\
Lemma~\ref{lem:jones49-setting} & \lean{HPC.Jones.jones49_setting}, \lean{HPC.Jones.jones49_simple_domains}. Jones: ``Because $M$ is homogeneous and contains a non-separating point, no point of $M$ separates $M$. Since $M$ is aposyndetic, $M$ is semi-locally-connected. Hence each element of $G$ is a simple domain.''\\
Lemma~\ref{lem:local-sep} & \lean{HPC.Jones.isLocalSepPoint_of_open_sub_tscc}, \lean{HPC.Jones.isLocalSepPoint_of_separation}.\\
Lemma~\ref{lem:frontier-of-jordan} & \lean{HPC.Jones.frontier_cc_of_jordan}.\\
Proposition~\ref{prop:case1} & \lean{HPC.Jones.exists_frontier_through}, \lean{HPC.Jones.countable_complDomains}, \lean{HPC.Jones.jones49_case1_core}, \lean{HPC.Jones.jones49_case1}. Jones ends with: ``Since no point of $M$ is of order~2 in $M$, each point of $J_i$ is a limit point of $M-J_i$. This contradicts a well known theorem (Baire).''\\
Lemma~\ref{lem:theta-sep} & \lean{HPC.Jones.theta_separation}, \lean{HPC.Jones.theta_separation_core}, with \lean{HPC.Jones.mem_cc_of_not_mem_domain} and \lean{HPC.Jones.side_eq_cc}.\\
Proposition~\ref{prop:case2} & \lean{HPC.Jones.exists_domain_meeting}, \lean{HPC.Jones.exists_theta_arcs}, \lean{HPC.Jones.jones49_case2_core}, \lean{HPC.Jones.jones49_case2}; the polygonal arc is \lean{HPC.Jones.exists_tarc_of_isOpen_plane}. Jones: ``It follows that $M-(A+B)$ is \emph{not} connected, but $M-B$ \emph{is} connected. Hence $A$ is a local separating point of $M$.''\\
Theorem~\ref{thm:jones49-1} & \lean{HPC.Jones.jones49_theorem1}, of type \lean{HPC.Jones.Jones49Thm1}.\\
Corollary~\ref{cor:jones49-2} & \lean{HPC.Jones.jones49_theorem2}, from \lean{HPC.Jones.theorem18}; the one-point space is again a counterexample to the printed statement, \lean{HPC.Jones.not_jones49_theorem2_point}.\\[4pt]
\multicolumn{2}{@{}l}{\emph{Chapter~\ref{ch:decomposable}}}\\
Definition~\ref{def:gprime} & \lean{HPC.Jones.gPrime}.\\
Proposition~\ref{prop:gprime} & \lean{HPC.Jones.isClosedMap_mk_gPrime}, \lean{HPC.Jones.mooreDecomp_gPrime}, \lean{HPC.Jones.jones55_plane}, \lean{HPC.Jones.jones55_embedding_data}, \lean{HPC.Jones.jones55_embedding}. This is the construction of \vfive{Corollary~9.16}; Jones states the upper semicontinuity of $G'$ without argument.\\
Theorem~\ref{thm:jones55-2} & \lean{HPC.Jones.jones55_circle}, \lean{HPC.Jones.jones55_theorem2'} (and \lean{HPC.Jones.jones55_theorem2}, with Theorem~\ref{thm:jones49-1} as a hypothesis).\\
Corollary~\ref{cor:jones55-remarks} & \lean{HPC.Jones.jones55_two_domains'}, \lean{HPC.Jones.jones55_elements_indecomposable'}, \lean{HPC.Jones.jones55_three_domains'}; the plane Jordan curve theorem used is \lean{HPC.Jones.jordan_curve_theorem}. Jones derives (b) from his 1951 paper; the formal proof derives it from (a) as above.\\
Proposition~\ref{prop:elements-pseudoarc} & \lean{HPC.BJ59.isPseudoArc_Lset'}, from \lean{HPC.HO16.isPseudoArc_of_indec_homogeneous} and \lean{HPC.HO16.ho16_theorem20}.\\
Definition~\ref{def:atomic} & \lean{HPC.BJ59.IsAtomicMap}, \lean{HPC.BJ59.IsAtomicCircPAMap}. The terminology is that of Domokos and Prajs.\\
Lemma~\ref{lem:open-quotient} & \lean{HPC.BJ59.isOpenMap_quotient_mk_of_continuous}.\\
Theorem~\ref{thm:decomposable} & \lean{HPC.BJ59.circle_or_atomic_of_decomposable}.\\
Remark~\ref{rem:circle-like} & \lean{HPC.BJ59.circle_or_circPAMap_of_decomposable}, with the map structure \lean{HPC.BJ59.IsCircPAMap}.\\[4pt]
\end{longtable}


\end{document}
